\documentclass[output=paper,colorlinks,citecolor=brown]{langscibook}
\ChapterDOI{10.5281/zenodo.21237018}
\author{Claudio Iacobini\affiliation{Università degli Studi Roma Tre} and Maria Pina De Rosa\affiliation{Università degli Studi di Salerno}}
%\ORCIDs{}

\title{Diachronic trends in the formation of Italian deadjectival verbs}

\abstract{This chapter provides an overview of deadjectival verb formation in the history of the Italian language by outlining the productivity and comparing the semantics of the three available verbalising processes, namely, suffixation (\textit{chiaro} ‘light, clear’ > \textit{chiar-eggiare} ‘clear up’), conversion (> \emph{chiar-ire} ‘clarify’) and parasynthesis (> \mbox{\emph{s-chiar-ire}} ‘lighten’).
While in the early periods of the language history parasynthesis is the best attested process, from the mid-19th century on suffixation becomes predominant. Significantly, it is most commonly through suffixation that complex adjectives (and thus most of those of more recent formation) are verbalised. Like parasynthesis, conversion, despite being well attested in the early stages, has seen a marked decline in productivity, and is overall the least productive of the three processes. Variation in productivity over time is also observed when comparing different derivational patterns (e.g., different suffixes) within a single process.
The semantic analysis of the derived verbs confirms that when adjectives are turned into verbs, this is mostly to express a change of state (‘make/become \textsc{Adj}’), which is the semantic space for which all three processes potentially compete. States and activities\is{activity} (‘be/appear \textsc{Adj}’,  ‘behave in a \textsc{Adj}’), on the other hand, are only encoded by suffixation and conversion, with a prevalence of the former. Parasynthesis  is reserved for expressing change-of-state events, as well as for deriving the largest number of verbs expressing an inchoative meaning (‘become \textsc{Adj}’) in the active form, i.e., without the anticausative marker -\emph{si}.
Overall, this study allowed us to reconstruct the development of Italian deadjectival verbalisation over a wide time frame, while also gaining insights into the factors that might have influenced it and possibly contribute to shaping current trends as well.
}


\IfFileExists{../localcommands.tex}{
   \addbibresource{../localbibliography.bib}
   \usepackage{tabularx,multicol}
\usepackage{url}
\urlstyle{same}

 
\usepackage{langsci-optional}
\usepackage{langsci-lgr}
\usepackage{langsci-gb4e}
\usepackage[linguistics,edges]{forest}
\usepackage{tikz-qtree}
\usetikzlibrary{positioning}
\usetikzlibrary{patterns.meta}
\usepackage{multirow}
\usepackage{pgfplots}
\usepackage{setspace}
\usepackage{amsmath,stackengine}
% \usepackage{fontspec}
% \usepackage{tipa} % TIPA is banned
\usepackage{langsci-textipa}
\usepackage{langsci-branding}
\usepackage{longtable}
\usepackage{tabto}

\usepackage{enumitem}

%\usepackage{lscape} %Querformat chapter 4

   % \newcommand*{\orcid}{}


\makeatletter
\let\thetitle\@title
\let\theauthor\@author
\makeatother

\newcommand{\togglepaper}[1][0]{
   \bibliography{../localbibliography}
   \papernote{\scriptsize\normalfont
     \theauthor.
     \titleTemp.
     To appear in:
    Laura Grestenberger, Viktoria Reiter \& Melanie Malzahn (eds.).
    \emph{Deadjectival verb formation in Indo-European and beyond}.
     Berlin: Language Science Press. [preliminary page numbering]
   }
   \pagenumbering{roman}
   \setcounter{chapter}{#1}
   \addtocounter{chapter}{-1}
}

\newbool{bookcompile}
\booltrue{bookcompile}
\newcommand{\bookorchapter}[2]{\ifbool{bookcompile}{#1}{#2}}

%Sonderzeichen
%idg palatales k
\newcommand{\kinvbreve}{k\raisebox{0.6ex}{\hspace{-0.05em}\char"0311}}

\newcommand{\jac}{%
  \textit{\char"0237}% italic dotless j (ȷ)
  \hspace{-0.01em}% fine-tune horizontal spacing
  \raisebox{0.01ex}{\char"0301}% raised combining acute
}

%Baltic circumflex l
\newcommand\ltilded{%
\hspace{-0.2em}%
  \stackon[-1.5pt]{l}{\smash{\kern0.2em\~{}}}%
  \hspace{-0.2em}%
}

%Slav. wedge+ double acute
\newcommand{\HighH}[1]{%
    \stackon[-1.2ex]{#1}{\kern0.1em\H{}\kern-0.1em}%
}

%Lith dot tilde VICKY diese b.-sl. Sonderzeichen BRINGEN MICH UM!!
\newcommand{\tildedot}[1]{%
    \ensurestackMath{%
        \stackon[-0.27ex]{% Tilde layer 
            \stackon[-0.15ex]{% Dot layer 
                #1%
            }{\hspace{0.1em}\cdot\hspace{-0.13em}}%
        }{\hspace{0.1em}\text{\~{}}\hspace{-0.13em}%
    }%
}}

%Lith dot acute
\newcommand{\dotacute}[1]{%
    \ensurestackMath{%
        \stackon[-1.1ex]{% acute layer 
            \stackon[-0.15ex]{% Dot layer 
                #1%
            }{\hspace{0.1em}\cdot\hspace{-0.12em}}%
        }{\hspace{0.1em}\text{\'{}}\hspace{-0.12em}%
    }%
}}



% \setlength{\emergencystretch}{2pt} 

 \pgfplotsset{compat=1.18}


\defcitealias{AbB_1}{AbB 1}
\defcitealias{AbB_7}{AbB 7}
\defcitealias{ABL_1}{ABL 1}
\defcitealias{Vries1977}{AEW}
\defcitealias{ALEW}{ALEW}
\defcitealias{AnGr}{AnGr}
\defcitealias{ARM_1}{ARM 1}
\defcitealias{BT}{BT}
\defcitealias{black_concise_2000}{CDA}
\defcitealias{CHD}{CHD}
\defcitealias{CT28}{CT 28}
\defcitealias{CT34}{CT 34}
\defcitealias{CT39}{CT 39}
\defcitealias{Cleasby1874}{CV}
\defcitealias{Beekes2010}{EDG}
\defcitealias{eDiAna}{eDiAna}
\defcitealias{eDIL}{eDIL}
\defcitealias{EDL}{EDL}
\defcitealias{Kroonen2013}{EDPG}
\defcitealias{ESJS}{ESJS}
\defcitealias{EVP}{EVP}
\defcitealias{EWA1988}{EWA}
\defcitealias{EWAI}{EWAia I}
\defcitealias{EWAII}{EWAia II}
\defcitealias{EWGP}{EWGP}
\defcitealias{GDLI1961}{GDLI}
\defcitealias{GPC}{GPC}
\defcitealias{GRADIT1999}{GRADIT}
\defcitealias{HW}{HW}
\defcitealias{IEW}{IEW}
\defcitealias{KAH_2}{KAH 2}
\defcitealias{KAR_1}{KAR 1}
\defcitealias{KAR_2}{KAR 2}
\defcitealias{KEWA}{KEWA}
\defcitealias{Labat_TDP}{Labat TDP}
\defcitealias{lambert_BWL}{Lambert BWL}
\defcitealias{LEIA}{LEIA}
\defcitealias{Rix2001}{LIV\textsuperscript{2}}
\defcitealias{LIVaddenda}{LIV\textsuperscript{2} {A}ddenda}
\defcitealias{LKG1971}{LKG}
\defcitealias{LKZ}{LKŽ}
\defcitealias{LP}{LP}
\defcitealias{MhdGr}{MhdGr}
\defcitealias{MLLVG1959}{MLLVG}
\defcitealias{NIL}{NIL}
\defcitealias{OgnS}{OgnS}
\defcitealias{ONP}{ONP}
\defcitealias{RIMA_2}{RIMA 2}
\defcitealias{TIM_2}{TIM 2}
\defcitealias{WAP}{WAP}
\defcitealias{YOS_2}{YOS 2}






\newindex{wan}{wdx}{wnd}{Form index}
\newcommand{\iw}[2]{\index[wan]{#1!{#2}\xspace}}
\newcommand{\iwi}[3][]{\index[wan]{#2!{#3}\xspace#1}{#3}}
\newcommand{\iwit}[3][]{\rule{0pt}{0pt}#1\index[wan]{#2!{#3}}{#3}}


\newindex{van}{vdx}{vnd}{Form index}
\newindex{ban}{bdx}{bnd}{Book index}

% \newcommand{\iw}[3][]{\index[wan]{#2!\textit{#3}}}


% \providecommand{\iwi}[3][]{\iw{#3}\textit{#3}}


\newcommand{\indexnote}[1]{{\footnotesize\sffamily\raggedright\normalfont\color{gray}#1}}

   %% hyphenation points for line breaks
%% Normally, automatic hyphenation in LaTeX is very good
%% If a word is mis-hyphenated, add it to this file
%%
%% add information to TeX file before \begin{document} with:
%% %% hyphenation points for line breaks
%% Normally, automatic hyphenation in LaTeX is very good
%% If a word is mis-hyphenated, add it to this file
%%
%% add information to TeX file before \begin{document} with:
%% %% hyphenation points for line breaks
%% Normally, automatic hyphenation in LaTeX is very good
%% If a word is mis-hyphenated, add it to this file
%%
%% add information to TeX file before \begin{document} with:
%% \include{localhyphenation}
\hyphenation{
    par-a-digm non-de-adj-ect-iv-al -pre-sent -pre-sents for-ma-tion-al in-do-ger-ma-ni-sche in-do-ger-ma-ni-schen Rei-chert Rei-ter dia-chrony ety-mo-lo-gi-sches Ost-ger-ma-ni-schen ger-ma-ni-schen alt-in-di-schen mor-pho-lo-gi-scher Alt-indo-ari-sch-en lexico-graphy Schrij-ver
}
% aktuell in Bibliographie: stud-ies, histor-ical Soci-età
\hyphenation{
    par-a-digm non-de-adj-ect-iv-al -pre-sent -pre-sents for-ma-tion-al in-do-ger-ma-ni-sche in-do-ger-ma-ni-schen Rei-chert Rei-ter dia-chrony ety-mo-lo-gi-sches Ost-ger-ma-ni-schen ger-ma-ni-schen alt-in-di-schen mor-pho-lo-gi-scher Alt-indo-ari-sch-en lexico-graphy Schrij-ver
}
% aktuell in Bibliographie: stud-ies, histor-ical Soci-età
\hyphenation{
    par-a-digm non-de-adj-ect-iv-al -pre-sent -pre-sents for-ma-tion-al in-do-ger-ma-ni-sche in-do-ger-ma-ni-schen Rei-chert Rei-ter dia-chrony ety-mo-lo-gi-sches Ost-ger-ma-ni-schen ger-ma-ni-schen alt-in-di-schen mor-pho-lo-gi-scher Alt-indo-ari-sch-en lexico-graphy Schrij-ver
}
% aktuell in Bibliographie: stud-ies, histor-ical Soci-età
   \boolfalse{bookcompile}
   \togglepaper[10]%%chapternumber
}{}

\begin{document}
\maketitle

 
\section{Introduction}

This contribution explores the derivation of deadjectival verbs\is{deadjectival!verb} in the Italian language over a timeline spanning from the origins (13\textsuperscript{th} century) to the end of the 20\textsuperscript{th} century. As Italian can avail itself of three different derivational processes to form verbs from adjectives -- namely \isi{suffixation}, \isi{conversion}, and \isi{parasynthesis} -- which exhibit overlapping semantic functions, our goal is to investigate how these processes have interacted and competed\is{competition} over time.

The study, based on a dataset consisting of all the attested verbs derived from the most frequent Italian adjectives  throughout history, outlines and compares the \isi{productivity} of the different verbalising processes by tracking novel coinages, while also taking into account the \isi{competition} between synonymous verbs, a pivotal factor contributing to the marginalisation and eventual obsolescence of lexemes. Indeed, it is by considering the relationship between new formations and the existing lexicon at a given point in time that more light can be shed on the dynamics determining the \isi{productivity} of word-formation processes.

The few studies on Romance verb formation that take a diachronic perspective point out an abundance of verbs formed from the same adjective in the first stages of the language and their reduction over time. This trend was first highlighted in the seminal research by \citet{Malkiel1941}, as well as in recent contributions on Portuguese and Spanish by, respectively, \citet{Rio-Torto2019} and \citet{Espejel2023}. The only diachronic investigation on Italian deadjectival verb\is{deadjectival!verb} formation available is \citet{GrossmannIacobini2025}, from which the present study differs in methodology, contributing new data to the documentation and analysis of the phenomenon.

This chapter is organised as follows. In Section \ref{sec:data}, we begin by describing the methodological choices underlying the data collection (\ref{sec:samples}) and the classification (\ref{sec:annotation}). In Section \ref{sec:results} we present the results. We first outline the diachronic distribution of verbs according to derivational process (\ref{sec:processes}) and pattern (\ref{sec:pattern}). Subsequently (\ref{sec:competition}), some insight into the phenomenon of \isi{competition} between the derivational processes is provided: Sections \ref{sec:adj1}--\ref{sec:adj2} focus on the adjectival side and highlight the relationship between the semantic and derivational features of the adjectives and their availability to serve as bases for verbal derivation,\is{derivation!verbal} while Section \ref{sec:adj3} is devoted to describing how the derivational processes have been employed for the expression of the various possible semantic types that are typically yielded in deadjectival verbalisation.\is{deadjectival!verb} Section \ref{sec:conclusion} wraps up the paper with a summary of the main findings.

\section{Data}\label{sec:data}
\subsection{Sample selection}\label{sec:samples}

The collection of verbs departed from a sample of 858 adjectives, i.e., those identified as the most frequent and widespread according to the \textit{Basic Vocabulary of the Italian Language} (\citealt{DeMauro1980}).\footnote{The adjective collection was carried out by consulting the BDVDB database (\citealt{ThorntonEtAl1997}), which contains the lemmas of the \textit{Basic Vocabulary of the Italian Language} annotated with orthographic, phonological, morphological, syntactic and etymological information, and provided with the date of first attestation. Adjectives from the \textit{Basic Vocabulary} were selected from those belonging to the categories \textit{Fundamental Lexicon} and \textit{High frequency}. The adjectives of the \textit{Basic Vocabulary} cover about 90\% of the adjectives that can be found in any modern Italian text. Besides, they are to a large extent of ancient attestation: More than 70\% are in fact attested as early as the 14\textsuperscript{th} century, and only 5\% are attested after the mid-19\textsuperscript{th} century.} These adjectives were selected to obtain a wide and representative list of deadjectival verbs\is{deadjectival!verb} formed over a broad historical range, extending from the 13\textsuperscript{th} century to the end of the 20\textsuperscript{th} century. A total of 917 verbs was gathered by means of a manual search in the \textit{Grande Dizionario della Lingua Italiana} (\citetalias{GDLI1961}) historical dictionary.

This method allowed us to identify the number and types of verbal derivatives\is{derivation!verbal} for each adjective, as well as the adjectives lacking documented verbal derivatives, thus offering a means to investigate the propensity of adjectives for verbalisation.

The use of the \citetalias{GDLI1961} was a necessary choice due to the lack of large diachronic corpora for Italian and of other electronically searchable resources. Our \isi{productivity} measures, based on dictionary data, quantify the number of verbs first documented in a given time period. The use of the \citetalias{GDLI1961} made it possible to compile a comprehensive list of verbs belonging to the diasystem of the Italian language. The verbs were subsequently sought in the \textit{Grande Dizionario della lingua Italiana} (\citetalias{GRADIT1999}), the largest dictionary of contemporary Italian and our source for the present-day usage of the verbs. Through this second search we could also identify the subset of verbs exclusively lemmatised in the historical dictionary that did not make it into the \citetalias{GRADIT1999}.

\subsection{Annotation criteria}
\label{sec:annotation}

\subsubsection{General}
The adjectives were classified based on morphological make-up following \citet{ThorntonEtAl1997}, who distinguish between simplex and complex adjectives, and identify, with regard to the latter category, the following adjective-forming processes: \isi{prefixation} (\iwi{It}{\textit{disonesto}} ‘dishonest’ < \iwi{It}{\textit{onesto}} ‘honest’), \isi{suffixation} (\iwi{It}{\textit{maschile}} ‘masculine’ < \iwi{It}{\textit{maschio}} ‘male’), \isi{compounding} (\iwi{It}{\textit{maleducato}} ‘ill-mannered’ < \iwi{It}{\textit{male}} ‘bad’ and \iwi{It}{\textit{educato}} ‘mannered’) and, finally, denominal \is{denominal!adjective} and deverbal \is{deverbal/deverbative!adjective} \isi{conversion} (\iwi{It}{\textit{marrone}} ‘brown’ < \iwi{It}{\textit{marrone}} ‘chestnut’, \iwi{It}{\textit{sveglio}} ‘awake, alert’ < \iwi{It}{\textit{svegliare}} ‘wake up’). Aiming at the maximal identification of morphological compositionality possible, they include within complex adjectives also those in which a morphological element is identifiable solely through the comparison of words sharing similar derivational characteristics, even if not fully analysable in discrete elements. The choice of adopting this classification is motivated by the diachronic setting of our study, and is further supported by finding that the distinction between synchronically fully-compositional and non-fully-compositional would not significantly impact the distribution of the dataset’s adjectives into morphologically complex and non-complex.\footnote{On the derivability of verbs from morphologically complex adjectives in Italian, see \citet{Grossmann1999}.}

In order to follow the evolution of the use of the different derivational processes and patterns over time, as well as to classify the verbs from a semantic and aspectual point of view, the verbs were annotated for a) first occurrence, b) current usage, c) verb-forming process and pattern expressing each process, and d) semantic and aspectual features. We discuss these criteria in more detail in the following sections.


 \subsubsection{First occurrence}     
The verbs’ date of first attestation was aligned to the temporal scan outlined in the MIDIA corpus (\citealt{DAchille2017}), which divides the history of the Italian language into five periods based on significant linguistic, literary, and cultural events:
    
\begin{enumerate}
        \item The period from the early 13\textsuperscript{th} century to 1375, spanning from the emergence of literature and vernacular writing in Tuscany to Boccaccio’s death and Coluccio Salutati’s appointment as Chancellor of Florence. \item From 1376 to 1532: This period includes the era of Humanism and the Renaissance, characterised by the evolution from “Silver Florentine” to the establishment of “Golden Florentine” as the literary norm, as advocated by Pietro Bembo in his \textit{Prose della volgar lingua} (1525). It concludes with the publication of the third edition of \textit{Orlando Furioso} by Ludovico Ariosto, embodying the realisation of Bembo’s linguistic principles in poetry. 
        \item From 1533 to 1691: The Late Renaissance, Mannerism and Baroque era, culminating in the publication of the third edition of the \textit{Vocabolario degli Accademici della Crusca} (1691), and the foundation of the Arcadia in 1690. 
        \item From 1692 to 1840: This epoch encompasses the Arcadia, the Enlightenment, and Romanticism. It culminates with the definitive edition of Alessandro Manzoni’s \textit{I Promessi Sposi}, which, drawing upon spoken Florentine, served as a linguistic model in post-unification Italy. 
        \item From 1841 to 1947: The era of the unification of Italy, World Wars, extending to the advent of the Italian Republic and the ratification of its constitution. \\Our periodisation includes a subsequent period corresponding to the Republican era. 
        \item From 1948 to the present: An epoch characterised by mass schooling, the spread of the mass media, and the resulting gradual affirmation in speech of the Italian standard language at the expense of the dialects (\citealt{DeMauro2014}).
    \end{enumerate}

 \subsubsection{Current usage}
 
 Verbs are grouped according to the usage labels of the \citetalias{GRADIT1999} dictionary. We distinguished the verbs that remain in present-day usage from those which have become obsolete. The latter group includes the verbs lemmatised in the \citetalias{GDLI1961} but not in the  \citetalias{GRADIT1999} (about 17\% of the original set) and those marked in the \citetalias{GRADIT1999} as obsolete and/or only attested in literary texts.


\subsubsection{Verb-forming process and pattern}

The processes available in Italian for the derivation of verbs from adjectival bases are \isi{suffixation}, \isi{conversion} and \isi{parasynthesis}. Each process is characterised by distinct patterns: Suffixation\is{suffixation} involves the addition of one of a set of suffixes, which all assign the verbs to the -\textit{are} inflectional class. Conversion \is{conversion} involves the use of inflectional marks of the -\textit{are} or -\textit{ire} class. Parasynthesis\is{parasynthesis} results from the combination of one of a set of prefixes\is{prefixation} and \isi{conversion}. 

Given the wide diachronic perspective adopted here, only the patterns allowing a comparison over the whole history of the language have been taken into consideration. The verbalising suffixes that have been used productively\is{productivity} from the origins until present day are \iwi{It}{-\textit{izzare}}, \iwi{It}{-\textit{eggiare}}, and \iwi{It}{-\textit{ificare}} (e.g., \iwi{It}{\textit{estremo}} ‘extreme’ > \iwi{It}{\textit{estremizzare}} ‘take to the extremes’, \iwi{It}{\textit{verde}} ‘green’ > \iwi{It}{\textit{verdeggiare}} ‘be/become green, lush’, \iwi{It}{\textit{intenso}} ‘intense’ > \iwi{It}{\textit{intensificare}} ‘intensify’). Verbs suffixed\is{suffixation} with \iwi{It}{-\textit{itare}} were also included because of the wide use of some of them (e.g., \iwi{It}{\textit{facile}} ‘easy’ > \iwi{It}{\textit{facilitare}}  ‘facilitate’, \iwi{It}{\textit{mobile}} ‘movable, moving’ > \iwi{It}{\textit{mobilitare}} ‘mobilise’) and their morphological and semantic transparency, despite their limited number and the current lack of \isi{productivity} of the suffix.

The suffixes \iwi{It}{-\textit{izzare}} and \iwi{It}{-\textit{eggiare}} go back to the same Greek etymon, -ιζειν.\iw{AG}{-ιζω, -ιζειν} The form \iwi{It}{-\textit{izzare}} comes to the Romance languages through learned borrowings from Latin, particularly used in religious and philosophical writing, while the suffix \iwi{It}{-\textit{eggiare}} goes back to the vernacular Latin form \iwi{Latin}{-\textit{idiare}}, manifesting clear Romance phonetic developments.

The suffix \iwi{It}{-\textit{ificare}} results from Latin compound verbs\is{compounding} like \textit{certificare} \iw{Latin}{\textit{certificāre}} (\iwi{Latin}{\textit{certum facere}}), \iw{Latin}{\textit{uīuificāre}} \textit{vīvificare}, \iwi{Latin}{\textit{mānsuēfacere}}, \iwi{Latin}{\textit{stupefacere}} and from other verbs derived either from nouns ending in \iwi{Latin}{-\textit{fex}} or from adjectives ending in \iwi{Latin}{-\textit{ficus}}. Suffixed\is{suffixation} verbs have only ever belonged to the default inflectional class (\iwi{Latin}{-\textit{āre}}).

Conversion \is{conversion} can be defined, following \citet[322]{Valera2015}, as “a word-formation process where the form of the converted item does not change, while its inflectional potential, its syntactic function and its meaning do, such that the item displays inflectional, syntactic and semantic properties of a new word class.” While well suited to languages with poor inflectional systems, this definition does not fully apply to derivation by \isi{conversion} in languages, like Italian, in which the inflectional features of verbs are marked by a rich set of inflectional endings, which clearly distinguish the categories of verbs from those of nouns and adjectives.

Conversion \is{conversion} assigns a verb to the -\textit{are}, by large the most productive,\is{productivity} or the \nobreakdash-\textit{ire} inflectional class, e.g., \iwi{It}{\textit{libero}} ‘free’ > \iwi{It}{\textit{liberare}} ‘set free’, \iwi{It}{\textit{scuro}} ‘dark’ > \iwi{It}{\textit{scurire}} ‘darken’.

Parasynthesis\is{parasynthesis} can be defined as the simultaneous application of two derivational processes (\isi{prefixation} and \isi{conversion}), which results in the coinage of prefixed verbs\is{prefixation} directly derived from an adjectival (or nominal) base (see \citealt{Iacobini2021}). In the current synchronic stage neither the non-prefixed verb nor the prefixed\is{prefixation} adjective is an actual word. However, as already noted in \citet{Iacobini2004}, and corroborated by a number of examples from the present dataset, in prior stages of the language, parasynthetic\is{parasynthesis} and converted verbs \is{conversion} formed from the same adjective were attested in the same time period.

For the data collection, we sought the verbs with an \isi{ingressive} change-of-state \is{change of state (CoS)} meaning (‘make/become \textsc{Adj}’) that are derived by prefixes\is{prefixation} still currently productive,\is{productivity} i.e., \iwi{It}{\textit{ad}-}, \iwi{It}{\textit{in}-}, \iwi{It}{\textit{s}-}, \iwi{It}{\textit{de}-/\textit{di}-/\textit{dis}-}, \iwi{It}{\textit{ri}-/\textit{ra}-/\textit{rin}-}, e.g., \iwi{It}{\textit{dolce}} ‘sweet’ > \iwi{It}{\textit{addolcire}} ‘sweeten’, \iwi{It}{\textit{grasso}} ‘fat’ > \iwi{It}{\textit{ingrassare}} ‘gain weight; fatten’, \iwi{It}{\textit{bianco}} ‘white’ > \iwi{It}{\textit{sbiancare}} ‘whiten’, \iwi{It}{\textit{lucido}} ‘shiny, polished, lucid’ > \iwi{It}{\textit{delucidare}} ‘elucidate’, \iwi{It}{\textit{giovane}} ‘young’ > \iwi{It}{\textit{ringiovanire}} ‘look younger; make sb./st. look younger’.


\subsubsection{Semantic and aspectual features} 

The semantic and aspectual classification is based on the correlation between the meaning of the adjective and the predicative function\is{predicative!function} expressed by the verb. Verbs can fall, based on telicity and dynamicity, into the categories of change of state,\is{change of state (CoS)} \isi{activity} and state. Change-of-state (CoS)\is{change of state (CoS)} and activity verbs share the feature dynamic, whereas \isi{activity} and state verbs have the feature atelic in common. Based on argument structure, a further distinction can be made, within change-of-state verbs,\is{change of state (CoS)} into \isi{causative} and \isi{inchoative}. Activity\is{activity} (behaviour) and state (copulative) verbs, on the other hand, are only \isi{intransitive}. \tabref{tab:key1} shows the interplay between these defining criteria.


\begin{table}
\caption{Semantic and aspectual types of deadjectival verbs}
\label{tab:key1}
\begin{tabularx}{\textwidth}{lQQQ}
\lsptoprule
        & \multicolumn{2}{c}{{Dynamic}}                                                             & Static\\
\cmidrule{2-4}
        & {Telic}                                            & {{Atelic}}                           & Atelic\\
\cmidrule{2-4} 
        & {Change of state\is{change of state (CoS)}}        & {Activity}\is{activity}              & {State}\\
\cmidrule{2-4} 
{Tr.}   & {Causative\is{causative}}                          &                                      & \\
        & \iwit{It}{\textit{legalizzare}} ‘legalise’         &                                      & \\
        & \iwit{It}{\textit{chiarire}} ‘clarify’             &                                      & \\
        & \iwit{It}{\textit{rallegrare}} ‘cheer up’          &                                      & \\
\cmidrule{2-4} 
{Intr.} & {Inchoative\is{inchoative}}                        & {Behaviour}                          & {Copulative}\\
        & \mbox{\iwit{It}{\textit{impallidire}} ‘turn pale’}      & \mbox{\iwi{It}{\textit{americaneggiare}}} & \mbox{\iwi{It}{\textit{pianeggiare}}} ‘be flat’\\
        & \mbox{\iwi{It}{\textit{curiosare}}} ‘snoop around’      & ‘act American'                       & \\
\lspbottomrule
\end{tabularx}
\end{table}


In most cases the verbalisation of an adjective results in a CoS \is{change of state (CoS)} verb, i.e., the situation most commonly lexicalised by deadjectival verbs\is{deadjectival!verb} is the causation and/or the acquisition of the state or property denoted by the adjectival base. Such verbs can occur in \isi{causative} or \isi{inchoative} constructions.\footnote{In line with the terminology generally employed in the Italian word-formation literature, our use of the terms “\isi{causative}/\isi{inchoative}” does not presuppose a distinction to “\isi{factitive}/\isi{fientive}”: We use the former pair as general terms for \isi{transitive} and \isi{intransitive} change-of-state\is{change of state (CoS)}  deadjectival verbs.\is{deadjectival!verb}} In \isi{causative} constructions, the subject referent is the agent that causes the change of state\is{change of state (CoS)} of a patient/theme. With reference to the relationship between the adjective and the verb, the meaning of a \isi{causative} verb can generally be paraphrased as ‘make (more) \textsc{Adj}’, e.g., \iwi{It}{\textit{legale}} ‘legal’ > \iwi{It}{\textit{legalizzare}} ‘legalise’, \iwi{It}{\textit{chiaro}} ‘light’ > \iwi{It}{\textit{chiarire}} ‘clarify’, \iwi{It}{\textit{allegro}} ‘cheerful’ > \iwi{It}{\textit{rallegrare}} ‘cheer up’.

In \isi{inchoative} constructions, the subject referent takes on the role of the patient/theme affected by the CoS,\is{change of state (CoS)} regardless of the cause. The meaning of the verb can thus be paraphrased as ‘become (more) \textsc{Adj}’, e.g., \iwi{It}{\textit{pallido}} ‘pale’ > \iwi{It}{\textit{impallidire}} ‘turn pale’, ‘go white’. These verbs may exhibit a \isi{causative}-\isi{inchoative} alternation, which can either lack morphological marking (e.g., \textit{gli zuccheri ingrassano Giovanni} ‘sugars make John fat’/\textit{Giovanni ingrassa} \iw{It}{\textit{ingrassare}} ‘John is getting fat’) or, more commonly, be marked by the clitic pronoun \iwi{It}{-\textit{si}}. For some verbs, the clitic pronoun is obligatorily expressed (e.g., \iwi{It}{\textit{grande}} ‘big’ > \textit{ingrandire} ‘make bigger’/\textit{ingrandirsi} ‘grow bigger’),\iw{It}{\textit{ingrandire}/\textit{ingrandirsi}} while in others, its use is optional (e.g., \iwi{It}{\textit{zitto}} ‘silent’ > \iwi{It}{\textit{zittire}/\textit{zittirsi}} ‘go silent’). In contemporary Italian, the prevailing method of conveying the \isi{inchoative} sense involves the use of the pronominal form. However, in earlier stages of the language, numerous verbs were employed in \isi{inchoative} constructions without the need for a clitic pronoun (cf.\ a.o.\ \citealt{Cennamo2015}; \citealt{Bentley2024}).

The CoS \is{change of state (CoS)} can manifest as either \isi{ingressive}, emphasising the transition into a state characterised by the possession \is{possessive} of the property denoted by the adjective from a preceding state, or \isi{egressive}, if the departure from the original state is highlighted during the transition from one state to another (cf.\ \citealt{Grossmann1994}; \citealt{Iacobini2021}). Put differently, the CoS\is{change of state (CoS)} may involve either the acquisition of a property not previously held or, conversely, the deprivation or loss of a property. In verbs where the adjective signifies a scalar, gradable attribute, and the scale is unbounded, the CoS \is{change of state (CoS)} occurs as an approximation or deviation towards an indefinite endpoint, representing a transition between relative states. This transition involves either an augmentation or diminution of a property already possessed to some degree.

Unlike denominal verbs,\is{denominal!verb} whose interpretation relies on context and on some kind of knowledge regarding the usual or expected use of the noun referent, deadjectival verbs\is{deadjectival!verb} do not normally require any encyclopaedic integration, as their essential function is to relate the property expressed by the adjective with the object referent in \isi{causative} constructions or with the subject in \isi{inchoative} constructions.

Deadjectival verbs \is{deadjectival!verb} also lexicalise situations in which no change of state \is{change of state (CoS)} is predicated. Within this category, further distinctions can be drawn, based on dynamicity and argument structure, into copulative and behaviour verbs. Copulative verbs denote an inherent or permanent mode of being of the subject referent, which can be paraphrased as ‘be \textsc{Adj}’, e.g., \iwi{It}{\textit{piano}} ‘flat’ > \iwi{It}{\textit{pianeggiare}} ‘be flat’.

Behaviour verbs are dynamic verbs paraphrasable as ‘behave in the manner denoted by \textsc{Adj}’ or ‘appear \textsc{Adj}’. The subject referent refers to an entity that demonstrates or possesses a certain property to varying degrees or intermittently over time, e.g., 
\iwi{It}{\textit{americano}} ‘American’ > \iwi{It}{\textit{americaneggiare}}  ‘play the yankee, act American’, \iwi{It}{\textit{curioso}} ‘curious’ > \iwi{It}{\textit{curiosare}} ‘nose around, snoop’, \iwi{It}{\textit{folle}} ‘deranged, lunatic’ > \iwi{It}{\textit{folleggiare}} ‘frolic, paint the town red’.

The analysis also revealed the presence of a small number of verbs that do not fit into the grid identified by the intersection of the parameters indicated so far. In this case, transitivity\is{transitive} does not coincide with a causation in the object referent of the property expressed by the base adjective, e.g., \iwi{It}{\textit{frequente}} ‘frequent’ > \iwi{It}{\textit{frequentare}} ‘attend, frequent, date’, \iwi{It}{\textit{violento}} ‘violent’ > \iwi{It}{\textit{violentare}} ‘rape’. These verbs lack the usual semantic transparency characterising the \isi{transitive} verbs previously mentioned, as the property expressed by the base adjective does not denote the final state or property of the object referent, but may be differently interpreted to refer, for example, to the subject referent, or to the event as a whole (cf.\ also \citealt[108--112]{Fabregas2023}). Due to their extremely small number, these verbs are not included in the quantitative evaluations in Section \ref{sec:adj2}.

\section{Results}\label{sec:results}
 \subsection{The derivational processes}\label{sec:processes}


\tabref{tab:key2} shows the percentage distribution of all verbs in the dataset by process and time period. Parasynthetic\is{parasynthesis} verbs are the most numerous, accounting for 45.6\% of the total. Suffixed\is{suffixation} verbs account for approximately a third of the total (32.7\%), whereas converted verbs \is{conversion} make up the smallest group (21.7\%). While about half of the verbs are formed during the initial two time periods, all three processes remain productive\is{productivity} until present day.

\begin{table}
\caption{Process productivity over time: Percentage of new coinages from the total of verbs}
\label{tab:key2}
\begin{tabularx}{\textwidth}{Xrrrrrrr}
\lsptoprule
 &  \textbf{I:}  &  \textbf{II:}  & \textbf{III:}  & \textbf{IV:}  & \textbf{V:} &  \textbf{VI:} & \\
  \textbf{Process} &   \textbf{1200–} &  \textbf{1376–} &  \textbf{1533–} &  \textbf{1693–} &  \textbf{1841–} &  \textbf{1947–} & \textbf{Total}\\
  \midrule
 Suffixation\is{suffixation} & {3.7} & {2.3} & {5.2} & {3.9} & {8.9} & {8.7} & \textbf{32.7}\\
Conversion\is{conversion} & {10} & {4.5} & {2.8} & {1.4} & {2.1} & {0.9} & { \textbf{21.7}}\\
{Parasynthesis\is{parasynthesis}} & {19.2} & {7.1} & {8.1} & {3.8} & {5.9} & {1.5} & {\textbf{45.6}}\\
\midrule
 \textbf{Total} & \textbf{32.9} & \textbf{13.9} & \textbf{16.1} &  \textbf{9.1}& \textbf{16.9} & \textbf{11.1} &  \textbf{100}\\
\lspbottomrule
\end{tabularx}
\end{table}

\tabref{tab:key3} shows the percentage of verbs formed by the different processes for each of the time periods. While \isi{parasynthesis} is the best attested process in each of the first three periods, from the end of the 17\textsuperscript{th} century it starts being overridden by \isi{suffixation}, which becomes the most prominent process from the mid-19\textsuperscript{th} century onwards. A remarkable decrease in \isi{productivity} concerns also \isi{conversion}, which, while being well attested in the first two periods, shows a significant decrease in use from the middle of the 16\textsuperscript{th} century onwards.


\begin{table}
\caption{Process productivity over time: Percentage of new coinages per time period}
\label{tab:key3}
\begin{tabularx}{\textwidth}{Xrrrrrr}
\lsptoprule
 &  \textbf{I:}  &  \textbf{II:}  &  \textbf{III:}  & \textbf{IV:}  & \textbf{V:} &  \textbf{VI:}  \\
  \textbf{Process} &   \textbf{1200–} &  \textbf{1376–} &  \textbf{1533–} &  \textbf{1693–} &  \textbf{1841–} &  \textbf{1947–} \\
  \midrule
{Suffixation\is{suffixation}} & {11.3} & {16.5} & {32} & {42.8} & {52.9} & {78.5}\\
{Conversion\is{conversion}} & {30.4} & {32.3} & {17.7} & {15.5} & {12.3} & {7.8}\\
{Parasynthesis\is{parasynthesis}} & {58.3} & {51.2} & {50.3} & {41.7} & {34.8} & {13.7}\\
\midrule
{\textbf{Total}}  & \textbf{100} & \textbf{100} & \textbf{100} & \textbf{100} & \textbf{100} & \textbf{100}\\
\lspbottomrule
\end{tabularx}
\end{table}

Being aware that selecting high-frequency adjectives might result in an over-representation of the processes employed in the earliest stages, we compared the distribution of the verbs collected with the distribution of all derived verbs lemmatised in the \citetalias{GRADIT1999} (represented, respectively, by the solid and the dotted line in \figref{fig:key:1}), with the aim of checking whether our collection method significantly skewed the proportions of the verbal coinages across the different time periods.

\begin{figure}
\includegraphics[width=0.8\textwidth]{IDR1}
\caption{\label{fig:key:1} Coinage over time of all complex verbs (deadjectival, denominal, deverbal) vs.\ deadjectival verbs collected in this study}
\end{figure}

Although the two distributions diverge in the earliest and latest time periods, they show a largely similar trend. As expected, the verbs in our dataset are slightly more numerous than those in the \citetalias{GRADIT1999} during the earliest period, whereas the reverse is true in the most recent one. Nonetheless, our data distribution does not appear significantly skewed, in that even in the two most recent time periods it closely mirrors that resulted from the \citetalias{GRADIT1999}, and also aligns with the findings of quantitative studies on 20\textsuperscript{th} century Italian word formation (\citealt{IacobiniThornton1992}, \citealt{PešaMatracki2006}, \citealt{Micheli2024}), which show a decrease in parasynthetic\is{parasynthesis} verbs and an increase in suffixed\is{suffixation} verbs when compared to previous stages of the language.

 \subsection{The derivational patterns} \label{sec:pattern}
\subsubsection{Suffixation} 

The formation of suffixed\is{suffixation} verbs is especially concentrated in the two most recent time periods, as shown in \tabref{tab:key:4}.


\begin{table}
\caption{Productivity of suffixes over time: Percentage of new coinages on the total of suffixed verbs}
\label{tab:key:4}
\begin{tabularx}{\textwidth}{Xrrrrrrr}
\lsptoprule
 &  \textbf{I:}  &  \textbf{II:}  &  \textbf{III:}  & \textbf{IV:}  & \textbf{V:} &  \textbf{VI:}  & \\
  \textbf{Pattern} &   \textbf{1200–} &  \textbf{1376–} &  \textbf{1533–} &  \textbf{1693–} &  \textbf{1841–} &  \textbf{1947–} & \textbf{Total}\\
  \midrule
\textit{-izzare} & 1 & 1.7 & 6.3 &  8.3 & 19.3 & {22.7} & {\textbf{59.3}}\\
\textit{-eggiare} & 5.3 & 3.7 & 8 & 2.3 & 5.3 & {2.7} & {\textbf{27.3}}\\
\textit{-ificare} & 4.7 & 1 & 0.7 & 1 & 2.7 & {1.3} & {\textbf{11.4}}\\
{\textit{-itare}} & {0.3} & {0.7} & {0.7} & {0.3} & {0} & {0} & {\textbf{2}}\\
\midrule
{\textbf{Total}} & {\textbf{11.3}} & {\textbf{7.1}} & {\textbf{15.7}} & {\textbf{11.9}} & {\textbf{27.3}} & {\textbf{26.7}} & {\textbf{100}}\\
\lspbottomrule
\end{tabularx}
\end{table}

\begin{table}
\caption{Productivity of suffixes over time: Percentage of new coinages per time period}\label{tab:key:5}
\begin{tabularx}{\textwidth}{Xrrrrrr}
\lsptoprule
 &  \textbf{I:}  &  \textbf{II:}  &  \textbf{III:}  & \textbf{IV:}  & \textbf{V:} &  \textbf{VI:}  \\
  \textbf{Pattern} &   \textbf{1200–} &  \textbf{1376–} &  \textbf{1533–} &  \textbf{1693–} &  \textbf{1841–} &  \textbf{1947–} \\
  \midrule
{\iwit{It}{-\textit{izzare}}} & {8.8} & {23.8} & {40.3} & {69.4} & {70.7} & {85}\\
{\iwit{It}{-\textit{eggiare}}} & {47.1} & {52.4} & {51.1} & {19.5} & {19.5} & {10}\\
{\iwit{It}{-\textit{ificare}}} & {41.2} & {14.3} & {4.3} & {8.3} & 9.8 & {5}\\
{\iwit{It}{-\textit{itare}}} & {2.9} & {9.5} & {4.3} & {2.8} & {0} & {0}\\
\midrule
{\textbf{Total}}  & \textbf{100} & \textbf{100} & \textbf{100} & \textbf{100} & \textbf{100} & \textbf{100}\\
\lspbottomrule
\end{tabularx}
\end{table}


\tabref{tab:key:5} shows the percentage of verbs formed by different suffixes in each period. In the earliest centuries, \iwi{It}{-\textit{eggiare}} is the most frequently used suffix, accounting for about half of suffixed\is{suffixation} verbs in each of the first three periods. From the end of the 17\textsuperscript{th} century, its \isi{productivity} decreases, becoming very scarce in contemporary Italian: From the second half of the 20\textsuperscript{th} century onwards, coinages in \iwi{It}{-\textit{eggiare}} constitute just over 10\% of suffixed deadjectival verbs.\is{suffixation} \is{deadjectival!verb} Conversely, from the mid-16\textsuperscript{th} century onwards, there is a significant increase in the \isi{productivity} of the suffix \iwi{It}{-\textit{izzare}}, which becomes the most employed suffix by the end of the 17\textsuperscript{th} century; its \isi{productivity} continues to grow steadily, so that from the second half of the 20\textsuperscript{th} century onwards, verbs in \iwi{It}{-\textit{izzare}} constitute about 85\% of suffixed\is{suffixation} formations. The success of this suffix is linked to the influence of Modern Latin, and, more recently, to contact with French and English, cf. \citet{FallyGoryczka2021}; \citet{Radtke1994}.

Of the verbalising suffixes, \iwi{It}{-\textit{ificare}} is the least productive,\is{productivity} forming 11.4\% of total suffixed\is{suffixation} verbs, and appears to be employed mainly in the first two temporal phases. The recent reduction of \isi{productivity} that is shown here may be overrepresented due to the characteristics of our dataset, which, being built from the highest-frequency adjectives, cannot fully account for the verbal coinages in the technical-scientific domains, where recourse to \iwi{It}{-\textit{ificare}} (as well as \iwi{It}{-\textit{izzare}}) is frequent.

The data concerning the \isi{productivity} of the three suffixes agrees with what is reported by \citet{Grossmann2004} for the second half of the 20\textsuperscript{th} century, according to which the vast majority of deadjectival suffixed\is{suffixation} verbs \is{deadjectival!verb} are in fact formed by means of \iwi{It}{-\textit{izzare}}, while recent \isi{productivity} of \iwi{It}{-\textit{ificare}} and \iwi{It}{-\textit{eggiare}} is rather low.

In contemporary Italian, the suffix \iwi{It}{-\textit{eggiare}} is predominantly used to lexicalise activities\is{activity} and states, while verbs formed in earlier stages more often also present CoS,\is{change of state (CoS)} \isi{causative} meanings. At present, these verbs (or, specifically, their \isi{causative} meanings, in the cases of polysemous verbs) are generally of low frequency or have fallen out of use, e.g., \iwi{It}{\textit{aspro}} ‘sour, harsh’ > \iwi{It}{\textit{aspreggiare}} ‘treat harshly’; \iwi{It}{\textit{falso}} ‘false’ > \iwi{It}{\textit{falseggiare}} ‘falsify’; \iwi{It}{\textit{stanco}} ‘tired’ > \iwi{It}{\textit{stancheggiare}} ‘tire’. Only in rare instances is the \isi{causative} meaning preserved in high frequency verbs, e.g., \iwi{It}{\textit{amaro}} ‘bitter, sour’ > \iwi{It}{\textit{amareggiare}} ‘embitter, sour’; \iwi{It}{\textit{pari}} ‘equal, even’ > \iwi{It}{\textit{pareggiare}} ‘level’, besides ‘tie, draw’.

Concerning the suffix \iwi{It}{-\textit{izzare}}, the opposite case is observed, whereby a number of verbs testifies to an earlier use of this suffix for the expression of an \isi{activity}, i.e., a manner of behaving. This is in fact the function originally exhibited in the early attestations in Latin, later reinterpreted as \isi{causative} (cf.\  \citealt{GibertSoteloPujolPayet2020}). Only rarely, however, do \iwi{It}{-\textit{izzare}} verbs exclusively express an \isi{activity} meaning, e.g., \iwi{It}{\textit{avaro}} ‘greedy, miserly’ > \iwi{It}{\textit{avarizzare}} ‘behave like a miser’, or \iwi{It}{\textit{ostile}} ‘hostile’ > \iwi{It}{\textit{ostilizzare}} ‘behave in a hostile way’. More often, both an \isi{activity} and a \isi{causative} meaning are documented for the same verb, e.g., \iwi{It}{\textit{ebreo}} ‘Jew/Jewish’ > \iwi{It}{\textit{ebraizzare}} ‘make Jewish or similar to Jews’ besides ‘imitate the Jews, follow Jewish customs and traditions’; \iwi{It}{\textit{unico}} ‘unique’ > \iwi{It}{\textit{unicizzare}} ‘be unique, behave as unique’ besides ‘unite, unify’. Such verbs are currently of low frequency or obsolete.

Due to this semantic overlap between the suffixes \iwi{It}{-\textit{izzare}} and \iwi{It}{-\textit{eggiare}}, especially observable in earlier periods (cf.\ also \citealt{Tronci2019}), various \isi{doublet}s are attested, most often lexicalising an \isi{activity}, e.g., the forms \iwi{It}{\textit{latineggiare}}/\mbox{\iwi{It}{\textit{latinizzare}}} (< \iwi{It}{\textit{latino}} ‘Latin’), both of which could express the meaning ‘make use of Latinisms’  in the past. In contemporary Italian, however, said meaning is predominantly (if not exclusively) expressed by the verb in \iwi{It}{-\textit{eggiare}}, while the verb in \iwi{It}{-\textit{izzare}} survives with the \isi{causative} meanings of ‘introduce to a people the Latin or neo-Latin language and culture’ and ‘give Latin form to words of another language’. The very few verbs in \iwi{It}{-\textit{izzare}} that currently lexicalise an \isi{activity} do not have an \iwi{It}{-\textit{eggiare}} counterpart, e.g., \iwi{It}{\textit{solidale}} ‘supportive, sympathethic’ > \iwi{It}{\textit{solidarizzare}} ‘sympathise (with sb.)’, \iwi{It}{\textit{sottile}} ‘subtle’ > \iwi{It}{\textit{sottilizzare}} ‘split hairs’. In the fewer cases of \isi{doublet}s in which the \iwi{It}{-\textit{eggiare}} verb used to express a \isi{causative} value, the verb has fallen out of use, with the rival in \iwi{It}{-\textit{izzare}} prevailing, e.g., \iwi{It}{\textit{ridicolo}} ‘ridicule’ > \iwi{It}{\textit{ridicoleggiare}}/\iwi{It}{\textit{ridicolizzare}} ‘ridicule’.

While specialisation over time is observed for \iwi{It}{-\textit{eggiare}} and \iwi{It}{-\textit{izzare}}, \iwi{It}{-\textit{ificare}} expresses from the beginning only a \isi{causative} value, accordingly to its etymological origin. It is therefore not surprising that in the dataset only one \isi{doublet} in \nobreakdash-\textit{ificare}\iw{It}{-\textit{ificare}} and \iwi{It}{-\textit{eggiare}} is attested, i.e., \iwi{It}{\textit{chiareggiare}}/\iwi{It}{\textit{chiarificare}} (< \iwi{It}{\textit{chiaro}} ‘light, clear’), among whose meanings the \isi{causative} ‘make clear, clarify’ is shared.  More numerous are \isi{doublet}s in \iwi{It}{-\textit{ificare}} and \iwi{It}{-\textit{izzare}}. Although the suffixes both contribute (primarily or exclusively) a \isi{causative} meaning, the verbs forming the \isi{doublet}s are never fully interchangeable, as in \iwi{It}{\textit{elettrizzare}}/\iwi{It}{\textit{elettrificare}} (< \iwi{It}{\textit{elettrico}} ‘electrical’), where the technical meaning of ‘charging a body with electricity’ is shared, but only \iwi{It}{\textit{elettrizzare}} is also used figuratively to mean ‘put in a state of excitement’. 

\subsubsection{Conversion} 

Converted verbs \is{conversion} are for the most part attested in the first two time periods, while being scarcely represented in the more recent ones (see \tabref{tab:key:6}). \tabref{tab:key:7} shows that the vast majority (83.4\%) of converted verbs \is{conversion} fall into the default inflectional class, the -\textit{are} class, and that only few verbs belong to the -\textit{ire} inflectional class: Most of them date back to the earliest periods, while from the 18\textsuperscript{th} century onwards only an extremely scarce number of -\textit{ire} verbs, of low-frequency or limited to literary register, is coined, e.g., \iwi{It}{\textit{imbecillire}} ‘become, make imbecile’ (1865), \iwi{It}{\textit{giallire}} ‘become yellow’ (1878). After the mid-20\textsuperscript{th} century no -\textit{ire} verbs are attested in our dataset.


\begin{table}
\caption{Productivity of conversion patterns over time: Percentage of new coinages on the total of verbs}
\label{tab:key:6}
\begin{tabularx}{\textwidth}{Xrrrrrrr}
\lsptoprule
 &  \textbf{I:}  &  \textbf{II:}  &  \textbf{III:}  & \textbf{IV:}  & \textbf{V:} &  \textbf{VI:}  & \\
  \textbf{Pattern} &   \textbf{1200–} &  \textbf{1376–} &  \textbf{1533–} &  \textbf{1693–} &  \textbf{1841–} &  \textbf{1947–} & \textbf{Total}\\
  \midrule
\iwit{It}{-\textit{are}}& 39.8 & 16.1 & 9 & 6.5 & 8 & {4} & {\textbf{83.4}}\\
\iwit{It}{-\textit{ire}} & 6.6 & 4.5 & 4 & 0 & 1.5 & {0} & {\textbf{16.6}}\\
\midrule
{\textbf{Total}} & {\textbf{46.2}} & {\textbf{20.6}} & {\textbf{13.1}} & {\textbf{6.5}} & {\textbf{9.5}} & {\textbf{4}} & {\textbf{100}}\\
\lspbottomrule
\end{tabularx}
\end{table}



\begin{table}
\caption{Productivity of conversion patterns over time: Percentage of new coinages per time period\label{tab:key:7}}
\begin{tabularx}{\textwidth}{Xrrrrrr}
\lsptoprule
 &  \textbf{I:}  &  \textbf{II:}  &  \textbf{III:}  & \textbf{IV:}  & \textbf{V:} &  \textbf{VI:}  \\
  \textbf{Pattern} &   \textbf{1200–} &  \textbf{1376–} &  \textbf{1533–} &  \textbf{1693–} &  \textbf{1841–} &  \textbf{1947–} \\
  \midrule
\iwit{It}{-\textit{are}} & 85.9 & 78 & 69.2 & 100 & 84.2 & 100\\
\iwit{It}{-\textit{ire}} & 14.1 & 22 & 30.8 & 0 & 15.8 & 0\\
\midrule
{\textbf{Total}} & {\textbf{100}} & {\textbf{100}} &  \textbf{100} & {\textbf{100}} & {\textbf{100}} & {\textbf{100}}\\
\lspbottomrule
\end{tabularx}
\end{table}

It is rather uncommon for the same adjective to yield more than one converted verb.\is{conversion} When this occurs, the adjective at hand is generally prolific and is the base for a number of verbs derived by other derivational processes, too. A case in point is \iwi{It}{\textit{chiaro}} ‘clear, light’ > \iwi{It}{\textit{chiarare}}/\iwi{It}{\textit{chiarire}}, both attested with the meanings ‘make clear; make intelligible, explain; to clear up’; for \iwi{It}{\textit{chiarare}} the meaning ‘illuminate’ is also attested. In present-day Italian only \iwi{It}{\textit{chiarire}} survives, with the meaning ‘clarify, explain’. Alongside these, the following suffixed\is{suffixation} and parenthetic verbs are attested: \iwi{It}{\textit{chiareggiare}} ‘clear up, clarify’, \iwi{It}{\textit{chiarificare}} ‘clarify’, \iwi{It}{\textit{dichiarare}} ‘declare’, \iwi{It}{\textit{dischiarare}} ‘clear up’, \iwi{It}{\textit{schiarare}} ‘illuminate, become clear’, \iwi{It}{\textit{schiarire}} ‘lighten, become clear, clear up’. Conversely, cases of adjectives being verbalised by both \isi{conversion} patterns and no other are rarely found. Among the very few examples is \iwi{It}{\textit{oscuro}} ‘dark’ > \iwi{It}{\textit{oscurare}} ‘make, become dark’ and \iwi{It}{\textit{oscurire}} ‘become dark’.

\subsubsection{Parasynthesis}\is{parasynthesis} 

Recourse to \isi{parasynthesis} is extensive in the early stages, with more than half (57.7\%) of all parasynthetic\is{parasynthesis} verbs coined in the first two periods. The contribution of this process to deadjectival verb\is{deadjectival!verb} formation decreases over time, becoming very limited in the most recent period (see \tabref{tab:key:8}). 


\begin{table}
\caption{Productivity of parasynthetic prefixes over time: Percentage of new coinages on the total of parasynthetic verbs}
\label{tab:key:8}
\begin{tabularx}{\textwidth}{Xrrrrrrr}
\lsptoprule
 &  \textbf{I:}  &  \textbf{II:}  &  \textbf{III:}  & \textbf{IV:}  & \textbf{V:} &  \textbf{VI:}  & \\
  \textbf{Prefix}\is{prefixation} &   \textbf{1200–} &  \textbf{1376–} &  \textbf{1533–} &  \textbf{1693–} &  \textbf{1841–} &  \textbf{1947–} & \textbf{Total}\\
  \midrule
{\iwit{It}{\textit{in}-}} & {12.5} & {6.2} & {9.9} & {3.3} & {5.5} & {1.9} & {\textbf{39.3}}\\
{\iwit{It}{\textit{ad}-}} & {17.9} & {2.9} &  3.8 & {1.5} & {1.4} & {0.2} & {\textbf{27.7}}\\
{\iwit{It}{\textit{ri}-/\textit{ra}-/\textit{rin}-}} & {6.7} & {3.1} & {2.4} & {2.7} & {3.8} & {0.7} & {\textbf{19.4}}\\
{\iwit{It}{\textit{s}-}} & {1.9} & {2.9} & {1.4} & {1} & {1.9} & {0.5} & {\textbf{9.6}}\\
{\iwit{It}{\textit{de}-/\textit{di}-/\textit{dis}-}} & {3.1} & {0.5} & {0.2} & {0} & {0.2} & {0} & {\textbf{4}}\\
\midrule
{\textbf{Total}}  & \textbf{42.1} & \textbf{15.6} & \textbf{17.7} & \textbf{8.5} & \textbf{12.8} & {\textbf{3.3}} & {\textbf{100}}\\
\lspbottomrule
\end{tabularx}
\end{table}

The best attested prefixes\is{prefixation} are \iwi{It}{\textit{in}-} (39.3\%), \iwi{It}{\textit{ad}-} (27.7\%), and \iwi{It}{\textit{ri}-/\textit{ra}-/\textit{rin}-} (19.4\%), while the presence of \iwi{It}{\textit{s}-} and \iwi{It}{\textit{de}-/\textit{di}-/\textit{dis}-} is rather limited (respectively 9.6\% and 4\%). As shown in \tabref{tab:key:9}, the \isi{productivity} of the prefix\is{prefixation} \iwi{It}{\textit{in}-} is consistently high – the highest, in fact, from the second time period onwards. In the earliest period the prefix\is{prefixation} \iwi{It}{\textit{ad}-} is the most robustly attested, but then undergoes a quite notable decrease in use, becoming the second least productive\is{productivity} prefix\is{prefixation} from the mid-17\textsuperscript{th} century onward. While already being well attested in the earlier stages, from the fourth period onwards the prefixes \iwi{It}{\textit{ri}-/\textit{ra}-/\textit{rin}-} emerge as the most employed after \iwi{It}{\textit{in}-}.




\begin{table}
\caption{Productivity of parasynthetic prefixes over time: Percentage of new coinages per time period}
\label{tab:key:9}
\begin{tabularx}{\textwidth}{Xrrrrrr}
\lsptoprule
 &  \textbf{I:}  &  \textbf{II:}  &  \textbf{III:}  & \textbf{IV:}  & \textbf{V:} &  \textbf{VI:}   \\
  \textbf{Prefix}\is{prefixation} &   \textbf{1200–} &  \textbf{1376–} &  \textbf{1533–} &  \textbf{1693–} &  \textbf{1841–} &  \textbf{1947–} \\
  \midrule
{\iwit{It}{\textit{in}-}} & {29.6} & {39.9} & {55.4} & {40.1} & {42.6} & {57.2}\\
{\iwit{It}{\textit{ad}-}} & {42.6} & {18.5} & {21.6} & {17.1} & {11.1} & {7.1}\\
{\iwit{It}{\textit{ri}-/\textit{ra}-/\textit{rin}-}} & {15.9} & {20} & {13.5} & {31.4} & {29.6} & {21.4}\\
{\iwit{It}{\textit{s}-}} & {4.5} & {18.5} & {8.1} & {11.4} & {14.8} & {14.3}\\
{\iwit{It}{\textit{de}-/\textit{di}-/\textit{dis}-}} & {7.4} & {3.1} & {1.4} & {0} & {1.9} & {0}\\
\midrule
{\textbf{Total}}  & \textbf{100} & \textbf{100} & \textbf{100} & \textbf{100} & \textbf{100} & \textbf{100}\\
\lspbottomrule
\end{tabularx}
\end{table}



As for the overall less productive\is{productivity} prefixes,\is{prefixation} \iwi{It}{\textit{s}-} shows a significant increase in use from the first to the second time period (from 4.5\% to 18.5\%), and from then on, a rather constant usage. On the other hand, verbs prefixed\is{prefixation} with \iwi{It}{\textit{de}-/\textit{di}-/\textit{dis}-} are predominantly attested in the earlier stages, with very scarce recent coinages – an exception being \iwi{It}{\textit{dimagrire}} ‘lose weight, slim down’, first attested in 1855 and a highly frequent verb. The current widespread use of the prefixes\is{prefixation} \iwi{It}{\textit{de}-/\textit{di}-/\textit{dis}-} mainly concerns the formation of deverbal verbs\is{deverbal/deverbative!verb} to which they contribute an \isi{egressive} (i.e., reversative or privative) meaning.


From \tabref{tab:key:10}, which illustrates the assignment of parasynthetic\is{parasynthesis} verbs into the two inflectional classes, we can observe a significant difference in both the overall proportions and the diachronic development compared to those outlined for converted verbs. While in each time period converted verbs \is{conversion} are predominantly assigned to the -\textit{are} class, the distribution of parasynthetic\is{parasynthesis} verbs between the two classes is more balanced (although still with a prevalence of -\textit{are} verbs). In particular, the first two time periods are characterized by a predominance of -\textit{are} verbs, the third period sees a shift toward the prevalence of -\textit{ire} verbs, while in the fourth period the \isi{productivity} of the two classes is comparable. The trend observed from the second half of the 19\textsuperscript{th} century onwards is opposite to that of the earlier periods, with the majority of parasynthetic\is{parasynthesis} verbs being assigned to the -\textit{ire} class. In this way, the non-default -\textit{ire} class, which has almost ceased to accommodate verbs formed by \isi{conversion} since the late 17\textsuperscript{th} century, retains its \isi{productivity}, e.g., \iwi{It}{\textit{curioso}} ‘curious’ > \iwi{It}{\textit{incuriosire}} (1887) ‘intrigue, arouse curiosity’; \iwi{It}{\textit{piatto}} ‘flat’ > \iwi{It}{\textit{appiattire}} (before 1920) ‘flatten’; \iwi{It}{\textit{nervoso}} ‘nervous, tense’ > \iwi{It}{\textit{innervosire}} (1950) ‘annoy, irritate’; \iwi{It}{\textit{volgare}} ‘vulgar, coarse’ > \iwi{It}{\textit{involgarire}} (1969) ‘vulgarise, coarsen’. 

\begin{table}
\caption{Distribution of parasynthetic patterns}
\label{tab:key:10}
\begin{tabularx}{\textwidth}{Xrrrrrrr}
\lsptoprule
 &  \textbf{I:}  &  \textbf{II:}  &  \textbf{III:}  & \textbf{IV:}  & \textbf{V:} &  \textbf{VI:}  & \\
  \textbf{Pattern} &   \textbf{1200–} &  \textbf{1376–} &  \textbf{1533–} &  \textbf{1693–} &  \textbf{1841–} &  \textbf{1947–} & \textbf{Total}\\
  \midrule
{\textit{prefix}\textsc{-adj}\textit{-are}} \iw{It}{-\textit{are}} & {27.9} & {8.9} & {7.4} &  4.3 & {5} & {1} & {\textbf{54.5}}\\
{\textit{prefix-}\textsc{Adj}\textit{-ire}} \iw{It}{-\textit{ire}} & {14.1} & {6.7} & {10.3} & {4.1} & {7.9} & {2.4} & {\textbf{45.5}}\\
\lspbottomrule
\end{tabularx}
\end{table}



Finally, \tabref{tab:key:11} shows the overall occurrence of parasynthetic\is{parasynthesis} prefixes\is{prefixation} over the \iwi{It}{-\textit{are}} and \iwi{It}{-\textit{ire}} inflectional classes. Verbs prefixed\is{prefixation} with \iwi{It}{\textit{in}-} are more often assigned the -\textit{ire} class, but are also numerous in the -\textit{are} class. Two-thirds of verbs prefixed\is{prefixation} with \iwi{It}{\textit{ad}-} belong to the -\textit{are} class. Verbs prefixed with \iwi{It}{\textit{ri}-/\textit{ra}-/\textit{rin}-} do not show a marked preference towards either of the two classes. Verbs prefixed\is{prefixation} with \isi{ingressive} \iwi{It}{\textit{s}-} predominantly fall within the \iwi{It}{-\textit{are}} class, but the percentage of \iwi{It}{-\textit{ire}} verbs is also high. Finally, the few verbs prefixed\is{prefixation} with \isi{ingressive} \iwi{It}{\textit{de}-/\textit{di}-/\textit{dis}-} almost exclusively belong to the \iwi{It}{-\textit{are}} class: In fact, no -\textit{ire} verb is attested with the prefixes\is{prefixation} \textit{dis}- and \textit{de}-, while by the pattern \textit{di-}\textsc{Adj}\textit{-ire} \iw{It}{-\textit{ire}} only the aforementioned \iwi{It}{\textit{dimagrire}} ‘slim down’ (< \iwi{It}{\textit{magro}} ‘slim’) is derived.


\begin{table}
\caption{Distribution of parasynthetic prefixes over the inflectional classes}
\label{tab:key:11}
\begin{tabularx}{\textwidth}{Xrrrrr}
\lsptoprule
{\textbf{Infl.\ class}}& \multicolumn{5}{c}{\textbf{Prefix}}\is{prefixation}\\
\midrule
 & \multicolumn{1}{l}{\iwit{It}{\textit{in}-}} & \multicolumn{1}{l}{\iwit{It}{\textit{ad}-}} & \multicolumn{1}{l}{\iwit{It}{\textit{ri}-/\textit{ra}-/\textit{rin}-}} & \multicolumn{1}{l}{\iwit{It}{\textit{s}-}} & \multicolumn{1}{l}{\iwit{It}{\textit{de}-/\textit{di}-/\textit{dis}-}}\\
\iwit{It}{-\textit{are}} & {41.5} & {66.4} & {54.4} & {57.5} & {94.1}\\
\iwit{It}{-\textit{ire}} & {58.5} & {33.6} & {45.6} & {42.5} & {5.9}\\
\midrule
{\textbf{Total}} & {\textbf{100}} & {\textbf{100}} & {\textbf{100}} & {\textbf{100}} & {\textbf{100}}\\
\lspbottomrule
\end{tabularx}
\end{table}



Parasynthesis\is{parasynthesis} is the process by which the most derivatives from the same adjective are formed. These are largely synonyms, e.g., \iwi{It}{\textit{caro}} ‘expensive, pricey’ > \iwi{It}{\textit{incarare}}/\iwi{It}{\textit{incarire}} ‘become more expensive’, both obsolete; \iwi{It}{\textit{largo}} ‘broad, large’ > \iwi{It}{\mbox{\textit{allargare}}}/\iwi{It}{\textit{slargare}} ‘enlarge’; \iwi{It}{\textit{ruvido}} ‘rough’ > \iwi{It}{\textit{arruvidare}} \mbox{(obsolete)/\mbox{\textit{irruvidire}}} ‘roughen’. Possible differences in meaning do not correlate systematically with the different patterns, e.g., \iwi{It}{\textit{profondo}} ‘deep’ > \iwi{It}{\textit{sprofondare}} ‘sink, collapse, make sb./st. collapse’/\textit{approfondire} ‘deepen, delve into’; \textit{sicuro} ‘sure, safe’ > \textit{assicurare} ‘assure, secure, fasten’/\iwi{It}{\textit{rassicurare}} ‘reassure’; \iwi{It}{\textit{rosso}} ‘red’ > \iwi{It}{\textit{arrossare}} ‘make red’/ \iwi{It}{\textit{arrossire}} ‘turn red, blush’.

 \subsection{Competition in deadjectival verb formation}
\label{sec:competition}

\subsubsection{Background}

The notion of \isi{competition} is at the centre of the current theoretical and methodological debate in morphology.\footnote{For recent comprehensive reviews on the theoretical issues and the methodological approaches developed for the study of \isi{competition} among affixes, and, more generally, in word formation, see \citet{HuygheVarvara2023}; \citet{NaganoEtAl2024}.} Italian deadjectival verb\is{deadjectival!verb} formation can serve as a useful testbed for the study of this phenomenon, in view of the several strategies available, both at the process and the pattern level, for producing the shift from the adjectival to the verbal lexical category.

Competition\is{competition} between derivational strategies often results in the instantiation of a word from a set of possible candidates, which, despite fulfilling all formal requirements, remain “virtual” (in the terminology of \citealt{Rainer2012}), as they get blocked by the already existing word. However, blocking is more of a tendency than a strict condition, and synonymous words can indeed be “actualised”, especially if we do not exclude partial synonymy as an instance of synonymy. This is the case with many adjectives from which multiple synonymous verbs have been derived. In such cases \isi{competition} becomes particularly manifest, but the interplay between derivational strategies can also be considered at a more abstract level, where these have overlapping preferences in base selection and can potentially fulfill the same semantic function.

To understand how \isi{competition} contributes to shape deadjectival verb\is{deadjectival!verb} formation, we examined the adjectives within our dataset, providing a quantitative overview of their propensity to verbalisation (\ref{sec:adj1}) and -- when this has taken place -- by which derivational strategies (\ref{sec:adj2}). We then explored how the different processes and patterns have overlapped or, conversely, divided up the labour in the expression of the main semantic and aspectual categories (\ref{sec:adj3}).

\subsubsection{Adjectives as bases for verbalisation: a quantitative overview}
\label{sec:adj1}

Out of the 858 adjectives examined, half have served as bases for verb derivation. Among the adjectives that have yielded verbal derivatives,\is{derivation!verbal} 48.8\% count only one, 24.2\% two, 12.6\% three, 14.4\% four or more.

An example of a very prolific base is \iwi{It}{\textit{molle}} ‘soft, mushy, weak’, which has yielded 11 verbs (listed, by period of coinage, in \tabref{tab:key:12}). All three processes have over time selected this adjective, from which two verbs are derived by \isi{conversion}, one by \isi{suffixation} (with \iwi{It}{-\textit{ificare}}), and as many as eight by \isi{parasynthesis}. As in many other cases, most verbs are formed in the earlier periods, but recent coinages are also attested. All derivatives lexicalise CoS:\is{change of state (CoS)} Some are attested only with a \isi{causative} meaning, i.e., ‘soften, appease, mollify’ (i.e., \mbox{\iwi{It}{\textit{mollificare}},} \mbox{\iwi{It}{\textit{ammollare}},} \iwi{It}{\textit{ammollire}}, \iwi{It}{\textit{immollare}}, \iwi{It}{\textit{smollare}}, \iwi{It}{\textit{mollire}}), others with both a \isi{causative} and an \isi{inchoative} one, i.e., ‘soften’ and ‘become soft, limp’ (\iwi{It}{\textit{mollare}}, \iwi{It}{\textit{immollire}}, \iwi{It}{\textit{rammollare}}, \iwi{It}{\textit{rammollire}}), and, finally, \iwi{It}{\textit{smollirsi}}, attested only in the pronominal form, is necessarily only used inchoatively.\is{inchoative} Of these derived verbs, six are obsolete or rare (\iwi{It}{\textit{mollificare}}, \iwi{It}{\textit{mollire}}, \iwi{It}{\textit{immollire}}, \iwi{It}{\textit{immollare}}, \iwi{It}{\textit{rammollare}}, \iwi{It}{\textit{smollirsi}}), while five are still in use and highly frequent (\iwi{It}{\textit{mollare}}, \iwi{It}{\textit{ammollare}}, \iwi{It}{\textit{ammollire}}, \iwi{It}{\textit{rammollire}}, \iwi{It}{\textit{smollare}}). While the verbs still used in present-day Italian share the general meaning ‘make soft, mollify’, they do vary in terms of usage context and collocations: \iwi{It}{\textit{rammollire}}, for example, is mostly used figuratively to mean ‘deprive of strength and energy’; \iwi{It}{\textit{ammollare}} is instead generally used to mean ‘dunk, soak’, i.e., ‘soften by immersing in a liquid’; finally, among the meanings of \iwi{It}{\textit{mollare}} (the most frequent and polysemous verb) are ‘release, leave, ditch, give, give in’, whereas ‘soften, weaken’ is marked in the \citetalias{GRADIT1999} as regional. No clear relation emerges between the time period of coinage and the survival of a verb. For example, among the most recently attested verbs, \iwi{It}{\textit{smollare}} survives, while \iwi{It}{\textit{smollirsi}} has not established itself in contemporary usage.

\begin{table}
\caption{\label{tab:key:12} Timeline of the coinages of derivatives from \emph{molle} ‘soft, mushy’}
\begin{tabularx}{\textwidth}{Xllllll}
\lsptoprule
 &  \textbf{I:}  &  \textbf{II:}  &  \textbf{III:}  & \textbf{IV:}  & \textbf{V:} &  \textbf{VI:}   \\
  \textbf{Process} &   \textbf{1200–} &  \textbf{1376–} &  \textbf{1533–} &  \textbf{1693–} &  \textbf{1841–} &  \textbf{1947–} \\
  \midrule
{Suffixation\is{suffixation}} & \iwit{It}{\textit{mollificare}} & & & & & \\
{Conversion\is{conversion}} & \iwit{It}{\textit{mollare}} &  & \iwit{It}{\textit{mollire}} &  &  & \\
{Parasyn.\is{parasynthesis}} &\iwi{It}{\textit{ammollare}} & \iwit{It}{\textit{immollire}} & & & \iwit{It}{\textit{smollare}}&\\
& \iwit{It}{\textit{ammollire}} & & & & \iwit{It}{\textit{smollirsi}}& \\
& \iwit{It}{\textit{immollare}} & & & & & \\
& \iwit{It}{\textit{rammollare}} & & & &    & \\
& \iwit{It}{\textit{rammollire}}  &  & &  &  & \\
\lspbottomrule
\end{tabularx}
\end{table}



A general overview of a single case such as this exemplifies the tendency to form multiple derivatives from the same base, often synonyms (whether full or partial), and to reduce their number over time. This trend is shared by other Romance languages and already documented by \citet{GrossmannIacobini2025} for Italian. This general tendency can be appreciated with reference to our data if we compare the ratio of the verbs derived from the adjectives first attested in the earliest time period to said adjectives (2,37\%) with the ratio of the verbs in present-day use to the total adjectives having ever yielded at least one verb (1,62\%). The reduction of the verbal derivatives\is{derivation!verbal} does not result, however, in a one-to-one verb-meaning correspondence.

In order to consider the complete landscape of \isi{competition} and the interaction between verbs derived from the same adjective, as well as the specialisation of meaning and the differentiation of semantic nuances in different domains or registers of verb usage, a detailed historical reconstruction of individual verbs is needed, but is beyond the scope and the dimension of this paper. An  analysis that is more limited in scope, but also more fine-grained  was carried out by \citet{Barbato2006}, who trace the history of the verbal derivatives\is{derivation!verbal} of the variants of the two adjectives \iwi{It}{\textit{bianco}} ‘white’ and \iwi{It}{\textit{chiaro}} ‘light, clear’ within the Italo-Romance varieties.

\subsubsection{Adjectives as bases for verbalisation: a morphosemantic analysis}\is{morphosemantics}
\label{sec:adj2}

Many factors possibly play a role in the instantiation and the establishment in the lexicon of a word derived by a certain strategy. Here we confine ourselves to observing the relevance, for verbalisation, of the adjectives’ morphological and semantic features. The adjectives that yield multiple verbs are, for the most part, morphologically simplex and express the prototypical property concepts\is{property concept (PC)} \citep{Dixon1982}, as shown in \tabref{tab:key:13}. Many of these properties can be understood both literally and figuratively.

\begin{table}
\caption{\label{tab:key:13}The adjectives with the highest number of verbal derivatives}
\begin{tabularx}{\textwidth}{p{2cm}X}
\lsptoprule
{\textsc{Colours}} & \mbox{\iwi{It}{\textit{azzurro}}} ‘blue’, \iwi{It}{\textit{bianco}} ‘white’, \iwi{It}{\textit{bruno}} ‘dark, dark brown’, \iwi{It}{\textit{nero}} ‘black’, \iwi{It}{\textit{pallido}} ‘pale’, \iwi{It}{\textit{rosa}} ‘pink’, \iwi{It}{\textit{rosso}} ‘red’, \iwi{It}{\textit{scuro}} ‘dark’, \iwi{It}{\textit{chiaro}} ‘light’\\
\tablevspace
{\textsc{Size}} & \mbox{\iwi{It}{\textit{basso}}} ‘short, low’, \iwi{It}{\textit{corto}} ‘short’, \iwi{It}{\textit{grande}} ‘big’, \iwi{It}{\textit{largo}} ‘wide’, \iwi{It}{\textit{lungo}} ‘long’, \iwi{It}{\textit{magro}} ‘skinny’, \iwi{It}{\textit{piccolo}} ‘small’, \iwi{It}{\textit{profondo}} ‘deep’, \iwi{It}{\textit{sottile}} ‘thin, subtle’\\
\tablevspace
\textsc{Physical} \textsc{properties} & \mbox{\iwi{It}{\textit{biondo}}} ‘blond’, \iwi{It}{\textit{calvo}} ‘bald’, \iwi{It}{\textit{curvo}} ‘curved’, \iwi{It}{\textit{denso}} ‘dense’, \iwi{It}{\textit{dolce}} ‘sweet’, \iwi{It}{\textit{doppio}} ‘double’, \iwi{It}{\textit{folto}} ‘thick’, \iwi{It}{\textit{freddo}} ‘cold’, \iwi{It}{\textit{lento}} ‘slow’, \iwi{It}{\textit{liquido}} ‘liquid’, \iwi{It}{\textit{lucido}} ‘shiny’, \iwi{It}{\textit{molle}} ‘soft, mushy’, \iwi{It}{\textit{morbido}} ‘soft’, \iwi{It}{\textit{nudo}} ‘bare, naked’, \iwi{It}{\textit{piano}} ‘flat’, \iwi{It}{\textit{secco}} ‘dry’, \iwi{It}{\textit{solido}} ‘solid’, \iwi{It}{\textit{sordo}} ‘deaf’, \iwi{It}{\textit{tenero}} ‘tender’, \iwi{It}{\textit{tiepido}} ‘lukewarm’, \iwi{It}{\textit{umido}} ‘moist’\\
\tablevspace
\textsc{Value} \mbox{\textsc{judgments}} & \mbox{\iwi{It}{\textit{bello}}} ‘nice, beautiful’, \iwi{It}{\textit{brutto}} ‘ugly’, \iwi{It}{\textit{buono}} ‘good’, \iwi{It}{\textit{caro}} ‘dear’, \iwi{It}{\textit{cattivo}} ‘bad’, \iwi{It}{\textit{cupo}} ‘gloomy’, \iwi{It}{\textit{grave}} ‘serious, deep’, \iwi{It}{\textit{misero}} ‘miserable, meager’, \iwi{It}{\textit{prezioso}} ‘precious’, \iwi{It}{\textit{puro}} ‘pure’, \iwi{It}{\textit{vero}} ‘true’, \iwi{It}{\textit{vile}} ‘vile’, \iwi{It}{\textit{vano}} ‘vain’, \iwi{It}{\textit{vivo}} ‘alive’\\
\tablevspace
\textsc{Human} \mbox{\textsc{attitudes}} & \mbox{\iwi{It}{\textit{curioso}}} ‘curious’, \iwi{It}{\textit{gentile}} ‘kind’, \iwi{It}{\textit{sereno}} ‘clear, peaceful’, \iwi{It}{\textit{superbo}} ‘arrogant’, \iwi{It}{\textit{triste}} ‘sad’\\
\lspbottomrule
\end{tabularx}
\end{table}

As stated above, the morphological structure of the adjectives notably impacts their propensity for verbal derivation.\is{derivation!verbal} 
As shown in \tabref{tab:key:14}, morphologically simplex adjectives, which account for just over a third of all adjectives, undergo verb derivation much more frequently than complex adjectives. In fact, 73.2\% of simplex adjectives are used to derive at least one verb. The percentage of verbs derived from complex adjectives is significantly lower: Only slightly more than one third (37.8\%) of complex adjectives has served as a base for verb derivation.



\begin{table}
\caption{\label{tab:key:14}Relationship between the morphological complexity of the adjectives and their propensity to be derived into verbs}
\begin{tabularx}{\textwidth}{Xrrrr}
\lsptoprule
\textbf{All adjectives} & & {\textbf{Verb-forming}} & {\textbf{Non verb-forming}} & {\textbf{Total}}\\
\midrule
Simplex & {34.8} & 73.2& 26.8 & \textbf{100}\\
Complex & {65.2} & 37.8 & 62.2 & \textbf{100}\\
\midrule
{\textbf{Total}} & {\textbf{100}} &  &  & \\
\lspbottomrule
\end{tabularx}
\end{table}


As shown in \tabref{tab:key:15}, the three verbalising processes show different propensities with regard to the selection of simplex versus complex adjectives.


\begin{table}
\caption{\label{tab:key:15}Relationship between the morphological complexity of the adjectival bases and the verb-forming processes}
\begin{tabularx}{\textwidth}{Xrrr}
\lsptoprule
{\textbf{Verbalising}} & \textbf{Verbs from}  & \textbf{Verbs from}  & \textbf{Verbs from} \\
\textbf{process}& \textbf{all \textsc{Adj}s} &\textbf{simplex \textsc{Adj}s} & \textbf{complex \textsc{Adj}s}\\
\midrule
Conversion\is{conversion} & 21.7 &23.4 & 19\\
Parasynthesis\is{parasynthesis} & 45.6 & 57.8 & 23.8\\
{Suffixation\is{suffixation}} & {32.7} & {18.8} & {57.2}\\
\midrule
{\textbf{Total}}  & {\bfseries 100} & {\bfseries 100} & {\bfseries 100}\\
\lspbottomrule
\end{tabularx}
\end{table}

The leftmost column recalls the percentage contribution of each process to the total pool of derived verbs. The middle column shows verbs derived from simplex adjectives, while the rightmost column shows the percentage of verbs derived from complex adjectives. We see that \isi{parasynthesis} and \isi{suffixation} are the two processes most influenced by the derivational structure of adjectives, with parasynthetic\is{parasynthesis} verbs being predominantly derived from simplex adjectives and suffixed\is{suffixation} verbs from complex ones. Converted verbs do not strongly relate with either category, and simplex bases are only slightly more numerous.\is{conversion}

The morphological make-up of the adjectives gains increased significance on the derivation of verbs when we correlate the adjective-forming process and the verb-forming one(s), as shown in \tabref{tab:key:16}. The two left-hand columns show the propensity of the different types of complex adjectives to be verbalised, while the subsequent three columns show the specific processes by which the verbs are derived.

\begin{table}
\caption{\label{tab:key:16}Relationship between the derivational structure of the adjectives and that of the resulting verbs}
\begin{tabularx}{\textwidth}{Xrrrrrr}
\lsptoprule
  {\textbf{Complex}}&  {\textbf{Non-verb}}& {\textbf{Verb}} & & \textbf{Conv.}& {\textbf{Parasyn.}\is{parasynthesis}} & {\textbf{Suffixed}} \\
 \textbf{\textsc{Adjs}} & \textbf{forming} & \textbf{forming} & {\textbf{Total}}&  \textbf{Verbs} & \textbf{Verbs} &  \textbf{Verbs}  \\
 \midrule
{Compound\is{compounding}} & 75 & {25} & {\bfseries 100} & 17 & {--} & 83 \\
{Prefixed\is{prefixation}} & {68} & {32} & {\bfseries 100} & {46} & {--} & {54} \\
{Suffixed\is{suffixation}} & {58} & {42} & {\bfseries 100} & {14} & {17} & {69} \\
{Converted\is{conversion}} & {72} & {28} & {\bfseries 100} & {42} & {26} & {32} \\
\lspbottomrule
\end{tabularx}
\end{table}

Parasynthetic\is{parasynthesis} verbs are only derived from suffixed\is{suffixation} and converted adjectives, whereas \isi{suffixation} and \isi{conversion} allow for verb derivation from all types of bases.

Verbalisation of adjectives formed by \isi{conversion} is limited by the fact that many of them are converted from verbs, which strongly reduces the possibility of further derivation into verbs, e.g., \iwi{It}{\textit{decidere}} ‘decide’ > \iwi{It}{\textit{deciso}} ‘decisive’ > \mbox{*\textit{decisare};} \iwi{It}{\textit{distendere}} ‘stretch, unwind’ > \iwi{It}{\textit{disteso}} ‘stretched, relaxed’ > *\textit{distesare}. The increasing use of \isi{suffixation} for verb derivation (cf.\ Section  \ref{sec:processes}) can be partly explained by the fact that suffixed\is{suffixation} adjectives lend themselves to further derivation through \isi{suffixation}. Given that new adjectives are predominantly coined through \isi{suffixation}, the pool of potential bases available for verb derivation via \isi{suffixation} is continuously growing.

Compound adjectives\is{compounding} tend to have only one derivative, mainly formed by \nobreakdash-\textit{izzare}\iw{It}{-\textit{izzare}} \isi{suffixation} (e.g., \iwi{It}{\textit{autonomo}} ‘autonomous’ > \iwi{It}{\textit{autonomizzare}} ‘autonomise’, \iwi{It}{\textit{democratico}} ‘democratic’ > \iwi{It}{\textit{democratizzare}} ‘democratise’), or by \iwi{It}{-\textit{are}} \isi{conversion}, e.g., \iwi{It}{\textit{uniforme}} ‘uniform’ > \iwi{It}{\textit{uniformare}} ‘uniform’.

Prefixed\is{prefixation} adjectives also tend to derive only a single verb. Exceptions with two derivatives also occur, e.g., \iwi{It}{\textit{impaziente}} ‘impatient’, which gives rise to two converted verbs,\is{conversion} \iwi{It}{\textit{impazientare}} and \iwi{It}{\textit{impazientire}}, both ‘become or make impatient’. Furthermore, from \iwi{It}{\textit{infinito}} ‘infinite’, both converted \iwi{It}{\textit{infinitare}} ‘make infinite’ and suffixed\is{suffixation} \iwi{It}{\textit{infinitizzarsi}} ‘become infinite’ are derived. The most commonly used suffix to form verbs from prefixed\is{prefixation} adjectives is \iwi{It}{-\textit{izzare}}, e.g., \iwi{It}{\textit{internazionale}} ‘international’ > \iwi{It}{\textit{internazionalizzare}} ‘internationalise’, \iwi{It}{\textit{immobile}} ‘immobile’ > \mbox{\iwi{It}{\textit{immobilizzare}}} ‘immobilise’.

Approximately half of the adjectives formed by \isi{conversion} form only one verb, e.g., \iwi{It}{\textit{tecnico}} ‘technical’ > \iwi{It}{\textit{tecnicizzare}} ‘technicalise’. Some adjectives converted from nouns, on the other hand, yield up to four derivatives, e.g., \iwi{It}{\textit{rosa}} ‘rose’ > \iwi{It}{\textit{rosa}} ‘pink’ > \iwi{It}{\textit{arrosare}}, \iwi{It}{\textit{inrosare}}, \iwi{It}{\textit{rosare}}, \iwi{It}{\textit{roseggiare}}, the former three meaning ‘colour, paint pink’, with \iwi{It}{\textit{rosare}} also figuratively meaning ‘illuminate with a pink light’, and \iwi{It}{\textit{roseggiare}} meaning ‘take on a pink hue due to the effect of a diffused pink light (of the sky) or due to the presence of numerous pink-coloured flowers (of a field)’.

Among complex adjectives, suffixed\is{suffixation} ones are those from which the highest number of verbs are derived, both in general and in terms of the number of derivatives per single adjective: Roughly one third of adjectives yields two or more derivatives: \iwi{It}{\textit{civile}} ‘civil’ > \iwi{It}{\textit{civilizzare}}, \iwi{It}{\textit{incivilire}} both ‘civilise’; \iwi{It}{\textit{cristiano}} ‘Christian’ > \iwi{It}{\textit{cristianificare}},  \iwi{It}{\textit{cristianizzare}} both ‘christianise’, \iwi{It}{\textit{cristianeggiare}} ‘behave as a Christian’; \iwi{It}{\textit{curioso}} ‘curious’ > \iwi{It}{\textit{curiosare}}, \iwi{It}{\textit{scuriosare}}, \iwi{It}{\textit{curioseggiare}} all ‘look around, snoop’, \iwi{It}{\textit{incuriosire}} ‘make curious, arouse curiosity’. 

The suffixes \iwi{It}{-\textit{ale}}, \iwi{It}{-\textit{bile}}, and \iwi{It}{-\textit{oso}} are among the most productive\is{productivity} for adjective formation. The verbs derived from adjectives in \iwi{It}{-\textit{ale}} are predominantly formed by \iwi{It}{-\textit{izzare}} (e.g., \iwi{It}{\textit{commerciale}} ‘commercial’ > \iwi{It}{\textit{commercializzare}} ‘commercialise’, \iwi{It}{\textit{nazionale}} ‘national’ > \iwi{It}{\textit{nazionalizzare}} ‘nationalise’), although verbs in \iwi{It}{-\textit{eggiare}} are also documented, e.g., \iwi{It}{\textit{orientale}} ‘Eastern’ > \iwi{It}{\textit{orientaleggiare}} ‘assume attitudes or characteristics typical of the East’. Adjectives in \iwi{It}{-\textit{bile}} also form verbs with \iwi{It}{-\textit{izzare}}, e.g., \iwi{It}{\textit{responsabile}} ‘responsible’ > \iwi{It}{\textit{responsabilizzare}} ‘make responsible’, \iwi{It}{\textit{sensibile}} ‘sensitive’ > \iwi{It}{\textit{sensibilizzare}} ‘sensitise, raise awareness’. The numerous adjectives in \iwi{It}{-\textit{oso}} tend not to be verbalised, with just over 10\% forming (mostly parasynthetic)\is{parasynthesis} verbs, e.g., \iwi{It}{\textit{pietoso}} ‘compassionate’ > \iwi{It}{\textit{impietosire}} ‘move to pity’, \iwi{It}{\textit{goloso}} ‘greedy’, ‘delicious’ > \iwi{It}{\textit{ingolosire}} ‘make sb's. mouth water’, \iwi{It}{\textit{nervoso}} ‘nervous’ > \iwi{It}{\textit{innervosire}} ‘irritate’.

\subsubsection{The verbal derivatives: the meanings expressed}
\label{sec:adj3}

Examining the semantic and aspectual properties of attested verbs can help to gain more insight into the \isi{competition} between processes and patterns for the same meanings. \tabref{tab:key:17} displays the percentage distribution of verbs over the main semantic and aspectual classes. Deadjectival verbs \is{deadjectival!verb} predominantly convey a CoS \is{change of state (CoS)} (86.5\% of total verbs, equivalent to the cumulative value of the first three columns), whereas verbs denoting states or activities\is{activity} make up 7.5\%. Additionally, alongside verbs falling into the two main classes, there exists a subset (6\%) of verbs that can express CoS \is{change of state (CoS)} along activities\is{activity} and \isi{stative} meanings. In their default, active form (i.e., non-pronominal), most of the CoS \is{change of state (CoS)} verbs are \isi{causative} (53.2\%), while about a quarter (24.8\%) can convey both \isi{causative} and \isi{inchoative} meanings. 8.5\% solely express an \isi{inchoative} value, of which a small subset is exclusively found in the pronominal form, i.e., is not attested in the default, active form.

\begin{table}
\caption{\label{tab:key:17}Distribution of verbs over the semantic and aspectual classes}
\begin{tabularx}{\textwidth}{Xrrrrr}
\lsptoprule
& {\textbf{CoS}}\is{change of state (CoS)} & & {\textbf{State/}} & \textbf{CoS}\textbf{ + state/} & \\
& & & \textbf{activity} & \textbf{activity}\is{activity} &{\textbf{Total}} \\
\midrule 
{Caus.\is{causative}} & {Inch.\is{inchoative}} & {Caus.\ + Inch.} &  &  &   \\
\midrule
{53.2} & {8.5} & {24.8} & {7.5} & {6} & {\bfseries 100}\\
\lspbottomrule
\end{tabularx}
\end{table}

\tabref{tab:key:18} shows the percentage contribution of each of the three processes in the expression of the different semantic types.

\begin{table}
\caption{\label{tab:key:18}Contribution of the derivational processes to the expression of the main semantic and aspectual types}
\begin{tabularx}{\textwidth}{Xrrrrr}
\lsptoprule
 & & {\textbf{CoS}}\is{change of state (CoS)}& & \textbf{State/} & {\textbf{CoS} \textbf{+}} {\textbf{state/}}  \\
  & & & & \textbf{{activity}}\is{activity} & {\textbf{activity}} \\
  \midrule
& {Caus.\is{causative}} & {Inch.\is{inchoative}} & {Caus.\ + Inch.} &  &   \\
\midrule
{Suffixation\is{suffixation}} & {39.2} & 6.7 & {4.1} & {66.7} & {64.1}\\
{Conversion\is{conversion}} & {23.2} & {16} & {16.5} & {27.3} & {30.2}\\
{Parasynthesis\is{parasynthesis}} & {37.6} & {77.3} & {79.4} & {6} & {5.7}\\
\midrule
{\textbf{Total}} & {\textbf{100}} & {\textbf{100}} & {\textbf{100}} & {\textbf{100}} & {\textbf{100}}\\
\lspbottomrule
\end{tabularx}
\end{table}

Regarding the expression of CoS, all three processes contribute in roughly equal proportions to the formation of \isi{causative} verbs, with \isi{conversion} slightly less represented than \isi{suffixation} and \isi{parasynthesis}. Conversely, verbs that exclusively express \isi{inchoative} meanings and verbs that occur in both \isi{inchoative} and \isi{causative} constructions are predominantly formed by \isi{parasynthesis}, with \isi{suffixation} playing a minor role. Parasynthesis\is{parasynthesis} is notably scarce among verbs expressing states and activities,\is{activity} as well as among those capable of denoting both CoS \is{change of state (CoS)} and states and activities.\is{activity} Suffixation\is{suffixation} emerges as the primary process for forming verbs denoting states and activities,\is{activity} while \isi{conversion} accounts for approximately 30\% of their derivation.

Tables 19 to 21 detail the expression of the main semantic types by the patterns within each process. \tabref{tab:key:19} displays the contribution of each suffix to the formation of CoS \is{change of state (CoS)} and of state and \isi{activity} verbs, from which a specialisation exhibited by individual suffixes emerges.



\begin{table}
\caption{\label{tab:key:19}Contribution of each suffix to the expression of the main semantic and aspectual types}
\begin{tabularx}{\textwidth}{Xrrrrrr}
\lsptoprule
 & & {\textbf{CoS}}\is{change of state (CoS)}& & \textbf{State/} & {\textbf{CoS} \textbf{+}} {\textbf{state/}} & \\
  & & & & \textbf{{activity}}\is{activity} & {\textbf{activity}} & \textbf{Total}\\
  \midrule
& {Caus.\is{causative}} & {Inch.\is{inchoative}} & {Caus.\ + Inch.} &  &  & \\
\midrule
{\iwit{It}{-\textit{izzare}}} & {48.2} & {1.3} & { 1.3} & { 1.1} & { 7.4} & { \textbf{59.3}}\\
{\iwit{It}{-\textit{eggiare}}} & { 1.3} & { 0} & { 1.1} & { 21} & { 4} & { \textbf{27.4}}\\
{\iwit{It}{-\textit{ificare}}} & { 10.4} & { 0.3} & { 0.6} & { 0} & { 0} & { \textbf{11.3}}\\
{\iwit{It}{-\textit{itare}}} & { 2} & { 0} & { 0} & { 0} & { 0} & { \textbf{2}}\\
\midrule
{\textbf{Total}} & { \textbf{61.9}} & { \textbf{1.6}} & { \textbf{3}} & { \textbf{22.1}} & { \textbf{11.4}} & { \textbf{100}}\\
\lspbottomrule
\end{tabularx}
\end{table}

The suffix \iwi{It}{-\textit{ificare}} is used exclusively in the formation of CoS \is{change of state (CoS)} verbs, whereas \iwi{It}{-\textit{izzare}}, besides being the most frequently employed suffix in the expression of CoS, is also exhibited by verbs denoting states and activities.\is{activity} The suffix \iwi{It}{-\textit{eggiare}} is unfrequently used in the derivation of CoS \is{change of state (CoS)} verbs, its main function being the expression of states and activities.\is{activity} The suffix \iwi{It}{-\textit{itare}} is found in only a very small number of \isi{causative} verbs.

Verbs formed by \isi{conversion} belonging to the \iwi{It}{-\textit{are}} class (83\%) significantly outnumber those in the \iwi{It}{-\textit{ire}} class (17\%). The contribution of each of the two classes to the expression of the different semantic and aspectual values can be seen in \tabref{tab:key:20}. We can see that among the verbs of the \iwi{It}{-\textit{are}} class there is a predominance of verbs with only \isi{causative} value, whereas the verbs of the \iwi{It}{-\textit{ire}} class can be mainly used in both \isi{causative} and \isi{inchoative} constructions. States and \isi{activity} verbs are almost exclusively of the -\textit{are} class.


\begin{table}
\caption{\label{tab:key:20}Percentage of converted verbs of the two classes in  the expression of the main semantic and aspectual types}
\begin{tabularx}{\textwidth}{lrrrrrr}
\lsptoprule
 & & {\textbf{CoS}}\is{change of state (CoS)} & & \textbf{State/} & {\textbf{CoS} \textbf{+}} {\textbf{state/}} & \\
  & & & & \textbf{{activity}}\is{activity} & {\textbf{activity}} & \textbf{Total}\\
  \midrule
& {Caus.\is{causative}} & {Inch.\is{inchoative}} & {Caus.\ + Inch.} &  &  & \\
\midrule
{\iwit{It}{-\textit{are}}} & { 59.7} & { 5.4} & { 11.4} & { 15.7} & { 7.8} & { \textbf{100}}\\
{\iwit{It}{-\textit{ire}}} & { 29.4} & { 8.8} & { 50.1} & { 2.9} & { 8.8} & { \textbf{100}}\\
\lspbottomrule
\end{tabularx}
\end{table}

The data regarding parasynthetic\is{parasynthesis} verbs are summarised in \tabref{tab:key:21}, wherein the various prefixes\is{prefixation} are grouped together due to their minor individual difference in contribution in this respect. This grouping also helps highlighting the differences in terms of inflectional class, and thus a more straightforward comparison with converted verbs.\is{conversion}


\begin{table}
\caption{\label{tab:key:21}Percentage of -\emph{are} and -\emph{ire} parasynthetic verbs in the expression of semantic and aspectual classes}
\begin{tabularx}{\textwidth}{Xrrrrrr}
\lsptoprule
 & & {\textbf{CoS}}\is{change of state (CoS)}& & \textbf{State/} & {\textbf{CoS} \textbf{+}}  & \\
  & & & & {\textbf{state/}}\textbf{{act.}} & {\textbf{act.}} & \textbf{Total}\\
  \midrule
& {Caus.\is{causative}} & {Inch.\is{inchoative}} & {Caus.\ + Inch.} &  &  & \\
\midrule
{Prefix\textit{\_are}}\iw{It}{-\textit{are}} & {32.9} & {6} & {13.2} & {1.4} & {0.7} & {\textbf{54.2}}\\
{Prefix\textit{\_ire}}\iw{It}{-\textit{ire}} & {9.4} & {8} & {28.4} & {0} & {0} & {\textbf{45.8}}\\
\midrule
{\textbf{Total}} & {\textbf{42.3}} & {\textbf{14}} & {\textbf{41.6}} & {\textbf{1.4}} & {\textbf{0.7}} & {\textbf{100}}\\
\lspbottomrule
\end{tabularx}
\end{table}

The data confirm the substantial percentage of CoS \is{change of state (CoS)} verbs, among which a significant proportion exhibits an \isi{inchoative} value and the near absence of verbs conveying states and activities.\is{activity} A difference that emerged between \isi{parasynthesis} and \isi{conversion} consists in a higher propensity of the former to derive \isi{inchoative} verbs -- among which particularly numerous are those belonging to the -\textit{ire} class. Some of the very few (only 6) verbs classified as states or activities\is{activity} may also be interpreted as prefixed\is{prefixation} verbs, cf. \iwi{It}{\textit{scuriosare}}/\iwi{It}{\textit{curiosare}} ‘snoop’, \iwi{It}{\textit{svigilare}}/\iwi{It}{\textit{vigilare}} ‘watch’. The only example of an unarguably parasynthetic\is{parasynthesis} verb is \iwi{It}{\textit{intristare}} ‘look sad’, not in common usage. Conversely, the verb \iwi{It}{\textit{impazzare}} is currently employed with an atelic, \isi{activity} meaning, ‘infuriate, manifest in a tumultuous way, rage, be all the rage’, while its \isi{inchoative} meaning, ‘go mad’, has fallen out of contemporary usage.

\section{Conclusions}
\label{sec:conclusion}

The study aimed to give a comprehensive overview of deadjectival verb\is{deadjectival!verb} formation in Italian, both in terms of the broad time span considered and the comparative approach used in describing the \isi{productivity} and semantics of the different verbalising strategies.

By taking the most frequent Italian adjectives as a starting point for the verb collection, we could distinguish the verb-forming from the non-verb-forming ones and identify the cases where multiple verbal derivatives\is{derivation!verbal} have been formed from one and the same base. There emerged that half of the 858 sampled adjectives have served as a base for verbal derivation,\is{derivation!verbal} of which 48.8\% have yielded a single derivative, 24.2\% two, 12.6\% three, and 14.4\% four or more.

Among the 917 verbs collected, parasynthetic\is{parasynthesis} derivatives are the highest in number (45.5\%), suffixed\is{suffixation} verbs account for about a third (32.6\%), while converted verbs make up the remaining 21.9\%. Most parasynthetic\is{parasynthesis} and converted verbs \is{conversion} are first attested in the earliest periods, while in 20\textsuperscript{th}-century Italian the vast majority of deadjectival verbs\is{deadjectival!verb} is instead derived by \isi{suffixation}, with a trend reversal beginning in the mid-19\textsuperscript{th} century.

The extensive use of \isi{parasynthesis} in the early centuries and of \isi{suffixation} in the later ones may, at least in part, be related to the processes’ preference for, respectively, morphologically simplex and complex bases, as revealed by the analysis of the adjectives’ morphological structure. The processes showing stronger preferences regarding the structure of the bases are \isi{parasynthesis} and \isi{suffixation}. Conversion,\is{conversion} on the other hand, lies in the middle, showing only a minimal preference for simplex adjectives. Overall, verbalisation involves simplex adjectives twice as often than complex ones. As for the latter, \isi{parasynthesis} is mostly employed to verbalise suffixed\is{suffixation} and converted adjectives, while \isi{suffixation} and \isi{conversion} show more flexibility in this regard. For compound\is{compounding} and suffixed\is{suffixation} adjectives \isi{suffixation} emerges as the preferred verbalising process.

In addition to examining base selection preferences, we also considered how the different processes and patterns are distributed with respect to the possible meanings expressed in deadjectival verbalisation.\is{deadjectival!verb} Indeed, an onomasiological perspective allows us to determine the extent to which the different strategies functionally overlap (and potentially compete),\is{competition} or differ.

The vast majority of deadjectival verbs,\is{deadjectival!verb} regardless of the derivational class, denote change-of-state\is{change of state (CoS)} events. The three processes contribute almost equally to the formation of \isi{causative} verbs, with \isi{suffixation} and \isi{parasynthesis} slightly more prevalent than \isi{conversion}, whereas verbs that exclusively express \isi{inchoative} meanings and those employed in both \isi{inchoative} and \isi{causative} constructions are mainly parasynthetic.\is{parasynthesis} While \isi{inchoative} meanings are normally achieved by adding the \isi{anticausative} particle \iwi{It}{-\textit{si}} to \isi{causative} verbs, a number of derivatives can convey such meaning in the active (i.e., non-pronominal) form: These are mostly parasynthetic\is{parasynthesis} and converted verbs \is{conversion} falling in the non-default, -\textit{ire} inflectional class. While change-of-state events\is{change of state (CoS)} can be expressed by all three processes, the encoding of states and activities\is{activity} is limited to \iwi{It}{-\textit{eggiare}} \isi{suffixation} and \iwi{It}{-\textit{are}} \isi{conversion}.

Considering the verbs’ first attestation and their persistence or exit from contemporary usage revealed two diachronic phenomena: A decrease in the number of verbs derived from the same adjective and an increasing presence of suffixed\is{suffixation} coinages. The expansion of \isi{suffixation} in verb formation is reasonably expected to continue, also in the light of the fact that \isi{suffixation} is the preferred process for the verbalisation of suffixed\is{suffixation} and compound\is{compounding} adjectives, which are, in turn, the adjectival classes most frequently enriched by new formations.

The original data and analyses presented in this research allowed for a wide-ranging diachronic investigation of the \isi{productivity} of the strategies involved in deadjectival verb\is{deadjectival!verb} formation and their functional distinctions. The insights provided contribute to a better understanding of the factors underlying the trend towards the decrease in the number of verbal lexemes over time and the current tendencies in the verbalisation of adjectives.

\section*{Abbreviations}
\begin{tabularx}{.45\textwidth}{lQ}
\textsc{Adj} & Adjective\\
caus. & causative\\
CoS & change of state\\
\end{tabularx}
\begin{tabularx}{.45\textwidth}{lQ}
inch. & inchoative\\
intr. & intransitive\\
tr. & transitive\\
\end{tabularx}

\sloppy\printbibliography[heading=subbibliography,notkeyword=this]
\end{document} 


